% =====================================================================
% BMED365 (Spring 2027) — Lab 4: Differentiable Neurosymbolic AI Report
% Annotated Source Code Monograph: kiwi_ns_lib.py (847 lines)
% Compile with: pdflatex kiwi_ns_lib_py.tex (run twice for TOC & LastPage)
% =====================================================================
\documentclass[10pt,a4paper]{article}

\usepackage[utf8]{inputenc}
\usepackage[T1]{fontenc}
\usepackage[english]{babel}
\usepackage{lmodern}
\usepackage{microtype}
\usepackage{graphicx}
\usepackage{booktabs}
\usepackage{xltabular}
\usepackage{tabularx}
\usepackage{array}
\usepackage[table]{xcolor}
\usepackage{amsmath,amssymb,mathtools}
\usepackage{enumitem}
\usepackage[margin=1.65cm,top=1.75cm,bottom=1.85cm,headheight=14pt]{geometry}
\usepackage{fancyhdr}
\usepackage{lastpage}
\usepackage{titlesec}
\usepackage[most]{tcolorbox}
\usepackage{listings}
\usepackage{float}
\usepackage{caption}
\usepackage[hyphens]{url}
\usepackage[
    colorlinks=true,
    linkcolor=uibblue,
    citecolor=uibblue,
    urlcolor=uibred,
    bookmarksnumbered=true,
    bookmarksopen=true,
    pdftitle={kiwi_ns_lib.py -- Annotated Shared PyTorch and Quantitative Imaging Library (BMED365 Spring 2027)},
    pdfauthor={Arvid Lundervold (University of Bergen & MMIV) -- BMED365 Spring 2027},
    pdfsubject={Annotated Source Code Monograph for Differentiable Neurosymbolic AI on 7T MRI}
]{hyperref}
\def\UrlBreaks{\do\/\do-\do_\do.}
% Allow clean line breaks after underscores in long filenames & Python identifiers
\let\origunderscore\_
\renewcommand{\_}{\origunderscore\allowbreak}

% =====================================================================
% COLORS (UiB / MMIV / BMED365 House Style)
% =====================================================================
\definecolor{uibblue}{RGB}{0, 61, 115}
\definecolor{uibred}{RGB}{175, 28, 44}
\definecolor{kiwigreen}{RGB}{27, 138, 90}
\definecolor{amberwarn}{RGB}{165, 95, 0}
\definecolor{amberbg}{RGB}{255, 250, 240}
\definecolor{lightgray}{RGB}{246, 248, 250}
\definecolor{rowalt}{RGB}{243, 247, 252}
\definecolor{codebg}{RGB}{245, 248, 252}
\definecolor{codehdr}{RGB}{224, 234, 246}
\definecolor{codekw}{RGB}{0, 75, 145}
\definecolor{codestr}{RGB}{155, 35, 45}
\definecolor{codecmt}{RGB}{90, 105, 118}
\definecolor{linenum}{RGB}{130, 142, 156}

% =====================================================================
% TYPOGRAPHY, HEADERS & ANNOTATION BOXES
% =====================================================================
\setlength{\parindent}{0pt}
\setlength{\parskip}{0.42em}
\setlength{\emergencystretch}{1.8em}
\setlist[itemize]{topsep=2pt,itemsep=2pt,leftmargin=1.45em}
\setlist[enumerate]{topsep=2pt,itemsep=2pt,leftmargin=1.65em}

\newcolumntype{L}[1]{>{\raggedright\arraybackslash}p{#1}}
\newcolumntype{C}[1]{>{\centering\arraybackslash}p{#1}}
\newcolumntype{R}[1]{>{\raggedleft\arraybackslash}p{#1}}
\newcolumntype{Y}{>{\raggedright\arraybackslash}X}

\titleformat{\section}{\large\bfseries\color{uibblue}}{\thesection}{0.65em}{}[\vspace{1pt}\color{uibblue!40}\titlerule]
\titleformat{\subsection}{\normalsize\bfseries\color{uibblue!88!black}}{\thesubsection}{0.55em}{}
\titleformat{\subsubsection}{\small\bfseries\color{uibblue!75!black}}{\thesubsubsection}{0.5em}{}
\titlespacing*{\section}{0pt}{1.15em}{0.42em}
\titlespacing*{\subsection}{0pt}{0.8em}{0.28em}
\titlespacing*{\subsubsection}{0pt}{0.55em}{0.2em}

\pagestyle{fancy}
\fancyhf{}
\fancyhead[L]{\footnotesize\textcolor{gray}{BMED365 (Spring 2027) $\cdot$ Annotated Library Reference}}
\fancyhead[R]{\footnotesize\textcolor{gray}{\texttt{kiwi\_ns\_lib.py} (847 lines)}}
\fancyfoot[L]{\footnotesize\textcolor{gray}{University of Bergen \& MMIV $\cdot$ \href{https://github.com/arvidl/BMED365-2027}{\textcolor{gray}{github.com/arvidl/BMED365-2027}}}}
\fancyfoot[R]{\footnotesize\textcolor{gray}{Page \thepage{} of \pageref{LastPage}}}
\renewcommand{\headrulewidth}{0.4pt}
\renewcommand{\footrulewidth}{0.2pt}

\fancypagestyle{firstpage}{
    \fancyhf{}
    \fancyfoot[L]{\footnotesize\textcolor{gray}{University of Bergen \& MMIV $\cdot$ \href{https://github.com/arvidl/BMED365-2027}{\textcolor{gray}{github.com/arvidl/BMED365-2027}}}}
    \fancyfoot[R]{\footnotesize\textcolor{gray}{Page \thepage{} of \pageref{LastPage}}}
    \renewcommand{\headrulewidth}{0pt}
    \renewcommand{\footrulewidth}{0.2pt}
}

% --- Boxes ---
\newtcolorbox{annotbox}[1][]{
    enhanced, breakable,
    colback=kiwigreen!5, colframe=kiwigreen!80!black,
    colbacktitle=kiwigreen!80!black, coltitle=white,
    fonttitle=\bfseries\small, title={#1},
    boxrule=0.55pt, arc=2.2pt, left=7pt, right=7pt, top=5pt, bottom=5pt
}

\newtcolorbox{designbox}[1][]{
    enhanced, breakable,
    colback=lightgray, colframe=uibblue,
    colbacktitle=uibblue, coltitle=white,
    fonttitle=\bfseries\small, title={#1},
    boxrule=0.6pt, arc=2.2pt, left=7pt, right=7pt, top=5pt, bottom=5pt
}

\newtcolorbox{warnbox}[1][]{
    enhanced, breakable,
    colback=amberbg, colframe=amberwarn,
    colbacktitle=amberwarn, coltitle=white,
    fonttitle=\bfseries\small, title={#1},
    boxrule=0.6pt, arc=2.2pt, left=7pt, right=7pt, top=5pt, bottom=5pt
}

\newtcolorbox{srcbox}[1][]{
    enhanced, breakable,
    colback=codebg, colframe=uibblue!35,
    colbacktitle=codehdr, coltitle=uibblue!90!black,
    fonttitle=\ttfamily\scriptsize\bfseries, title={#1},
    boxrule=0.55pt, arc=2.2pt, left=6pt, right=6pt, top=3pt, bottom=3pt,
    toptitle=2.5pt, bottomtitle=2.5pt
}

% --- Numbered Python Listing Style (exact original line numbers) ---
\lstdefinestyle{pyannot}{
    language=Python,
    basicstyle=\ttfamily\scriptsize,
    keywordstyle=\bfseries\color{codekw},
    stringstyle=\color{codestr},
    commentstyle=\itshape\color{codecmt},
    numbers=left,
    numberstyle=\tiny\ttfamily\color{linenum},
    numbersep=7pt,
    xleftmargin=20pt,
    xrightmargin=0pt,
    showstringspaces=false,
    breaklines=true,
    breakatwhitespace=false,
    breakindent=12pt,
    columns=fullflexible,
    keepspaces=true,
    upquote=true,
    aboveskip=2pt,
    belowskip=2pt,
    literate={±}{{$\pm$}}1
             {³}{{\textsuperscript{3}}}1
             {×}{{$\times$}}1
             {é}{{\'e}}1
             {ö}{{\"o}}1
             {Ł}{{\L{}}}1
             {Δ}{{$\Delta$}}1
             {α}{{$\alpha$}}1
             {–}{{--}}1
             {—}{{---}}1
             {≥}{{$\ge$}}1
             {≤}{{$\le$}}1
             {⊂}{{$\subset$}}1
             {∈}{{$\in$}}1
             {→}{{$\to$}}1
             {…}{{\ldots}}1
}

\newcommand{\ghblob}{https://github.com/arvidl/BMED365-2027/blob/main}
\newcommand{\repodir}{Lab4-Comp-Imag/Kiwi-challenge/docs/neurosymbolic-report}

% =====================================================================
\begin{document}
\thispagestyle{firstpage}

% --- TITLE BANNER ---
\begin{tcolorbox}[
    enhanced, colback=uibblue, colframe=uibblue,
    arc=3.5pt, left=12pt, right=12pt, top=9pt, bottom=9pt
]
{\Large\bfseries\color{white} \texttt{kiwi\_ns\_lib.py} --- Annotated Shared PyTorch \& Quantitative Imaging Library\\[2pt]
{\large Single Source of Truth for the Differentiable Neurosymbolic 7\,T MRI ``Thesis + Code'' Package}}\\[4pt]
{\normalsize\color{white!90} BMED365 --- Computational Imaging, Modeling and AI in Biomedicine (Spring 2027)}\\[3pt]
{\footnotesize\color{white!84}
Arvid Lundervold (University of Bergen \& MMIV) $\;\cdot\;$
847 lines $\;\cdot\;$ 1 dataclass, 5 \texttt{nn.Module} classes, 25 functions $\;\cdot\;$ 100\% pure PyTorch (CUDA / MPS / CPU)
}\\[6pt]
{\footnotesize
\href{\ghblob/\repodir/notebooks/kiwi_ns_lib.py}{\textcolor{white}{\underline{\texttt{kiwi\_ns\_lib.py} on GitHub}}}\hspace{12pt}
\href{\ghblob/\repodir/neurosymbolic-kiwi-report.pdf}{\textcolor{white}{\underline{Technical Report PDF}}}\hspace{12pt}
\href{\ghblob/\repodir/README.md}{\textcolor{white}{\underline{Package README}}}
}
\end{tcolorbox}

% --- ARCHITECTURE OVERVIEW ---
\begin{designbox}[Why This Annotated Reference Exists (Zero-Drift ``Thesis + Code'' Engineering)]
Every figure, ablation metric, calibration curve, spatial log-prior field, and whole-organ 3D NIfTI biomarker in \href{\ghblob/\repodir/neurosymbolic-kiwi-report.pdf}{\texttt{neurosymbolic-kiwi-report.pdf}} is generated by Companion Notebooks \texttt{01}--\texttt{05} importing this single 847-line module (or auto-fetching it verbatim from GitHub \texttt{main} in Google Colab). Centralizing all coordinate transforms, 3D architectures, fuzzy-logic losses, training loops, watershed/Mat\'ern decoders, and Hungarian evaluators here guarantees \textbf{zero code drift} between experiments and published LaTeX equations.

In the pages below, \textbf{every line of \texttt{kiwi\_ns\_lib.py} (L1--L847) is reproduced verbatim with its exact original line numbers}, interleaved with mathematical derivations, tensor-shape contracts, and biomedical software-engineering annotations.
\end{designbox}

\vspace{-0.2em}
{\small\tableofcontents}

\clearpage
% =====================================================================
\section*{Architectural Overview, Tensor Contracts \& API Cross-Reference}
\addcontentsline{toc}{section}{Architectural Overview, Tensor Contracts \& API Cross-Reference}
% =====================================================================

\subsection*{End-to-End 6-Stage Pipeline Architecture}
The library is organized into six cohesive layers matching the scientific workflow of the technical report:
\begin{enumerate}
    \item \textbf{Runtime \& Portability (L31--L86):} Hardware-agnostic device dispatch (\texttt{CUDA} $\to$ Apple \texttt{MPS} $\to$ \texttt{CPU}), async GPU stream synchronization for accurate wall-clock benchmarking, and multi-depth file discovery with automatic GitHub fallback for standalone Google Colab sessions.
    \item \textbf{Whole-Organ Geometry \& Biophysical Predicates (L89--L280):} Loads the $256 \times 256 \times 128$ 7\,T RARE MRI ($0.2344 \times 0.2344 \times 0.5000\,\text{mm}^3$) and curated $80 \times 80 \times 16$ Micro-GT (or a custom manually edited NIfTI via \texttt{gt\_nii\_override}), computes slice-normalized organ coordinates $(\tilde{r}, (\tilde{r}-0.43)^2, \tilde{p}_z)$, builds smooth differentiable predicates ($R_{\text{ring}}, B_{\text{non\_dark}}$), removes $B_1$ shading via 3D morphological closing, and constructs per-instance normalized 3D Euclidean Distance Transforms.
    \item \textbf{3D Neural \& Neurosymbolic Architectures (L282--L389):} Shares a single compact 3-level \texttt{Backbone3DUNet} ($\sim 340\text{k}$ params) across four controlled variants---pure neural \texttt{MultiTaskUNet3D} (\texttt{in\_ch=3}), naive \texttt{CoordConv3D\,+\,INorm} (\texttt{in\_ch=6}), explicitly factorized \texttt{AdditiveBayesOrganNet3D} ($+353$ params, un-normalized \texttt{Organ-GPS}), and nonlinear \texttt{CoupledOrganNet3D} ($+864$ params).
    \item \textbf{Differentiable Real Logic \& Multi-Task Losses (L391--L430):} Implements three canonical fuzzy implication families (Product/Reichenbach, \L{}ukasiewicz, G\"odel), LTN's generalized $p$-mean universal quantifier ($\varepsilon = 10^{-8}$ stabilizer), Soft Dice, and class-weighted multi-task BCE + Dice + distance MSE.
    \item \textbf{3D SSL Warm-Start \& Unified Training Harness (L432--L624):} Enforces strict spatial disjointness from the held-out Test Wedge across both 36 unlabeled SSL cubes and 40 polar-to-polar stratified unlabeled LTN blocks, supporting exact seed determinism, dual-LR \texttt{AdamW} + \texttt{CosineAnnealingLR}, 3D flip equivariance, and the causal columella hard-negative toggle (\texttt{use\_columella\_neg}).
    \item \textbf{Decoding, Steric Splitting, Evaluation \& Whole-Organ Inference (L626--L847):} Decodes probability + distance maps via physical-mm ellipsoidal 3D watershed, disentangles fused clumps via declarative Mat\'ern hard-core $K$-means ($V \ge 2.55\,\text{mm}^3$, $d_{\min} \ge 0.88\,\text{mm}$), evaluates exact Hungarian 3D Panoptic Quality ($\text{PQ} = \text{SQ} \times \text{RQ}$), calibrates $\tau_{\text{fg}}^*$ leakage-free on the training wedge, and tiles inference unmasked across all 128 axial slices.
\end{enumerate}

\subsection*{Neural Architecture Comparison \& Tensor-Shape Contract}
For any 3D input patch of spatial dimensions $(X, Y, Z)$ divisible by $4$ (e.g., training wedge $80 \times 40 \times 16$, unlabeled cubes $64 \times 40 \times 16$, whole-fruit tiles $80 \times 40 \times 16$), every model produces paired semantic logits $\mathbf{z}_{\text{sem}} \in \mathbb{R}^{B \times 1 \times X \times Y \times Z}$ and normalized center-distance predictions $\hat{\mathbf{d}} \in (0, 1)^{B \times 1 \times X \times Y \times Z}$:

{\scriptsize
\setlength{\tabcolsep}{4pt}
\renewcommand{\arraystretch}{1.20}
\rowcolors{2}{rowalt}{white}
\begin{tabularx}{\textwidth}{L{3.85cm} C{1.55cm} C{1.55cm} L{3.65cm} Y}
\toprule
\rowcolor{uibblue!12}\textbf{Class / Mode} & \textbf{Params} & \textbf{Delta} & \textbf{Inputs $(B, C, X, Y, Z)$} & \textbf{Coordinate Handling \& Auditability} \\
\midrule
\texttt{MultiTaskUNet3D(3)} (\texttt{baseline}) & 340,450 & --- & $\mathbf{X}_3$: $(B, 3, X, Y, Z)$ & Translation-equivariant CNN only; zero global organ awareness. \\
\texttt{MultiTaskUNet3D(6)} (\texttt{coordconv\_inorm}) & 341,746 & $+1{,}296$ & $[\mathbf{X}_3, \mathbf{g}]$: $(B, 6, X, Y, Z)$ & Naive concatenation into \texttt{InstanceNorm3d}; erases patch mean/scale (NB~02 proof). \\
\texttt{AdditiveBayesOrganNet3D} (\texttt{additive\_gps}) & \textbf{340,803} & \textbf{$+353$} & $\mathbf{X}_3$ $(B,3,\cdot)$ \& $\mathbf{g}$ $(B,3,\cdot)$ & Un-normalized $1{\times}1{\times}1$ MLP; explicit $z_{\text{app}} + z_{\text{GPS}}$; directly auditable via \texttt{gps\_log\_prior()}. \\
\texttt{CoupledOrganNet3D} (\texttt{coupled\_gps}) & 341,314 & $+864$ & $\mathbf{X}_3$ $(B,3,\cdot)$ \& $\mathbf{g}$ $(B,3,\cdot)$ & Un-normalized GPS embedding concatenated with decoder features; allows nonlinear interaction. \\
\bottomrule
\end{tabularx}
}

\subsection*{Complete Symbol Index \& Cross-Notebook Usage Matrix (All 31 Definitions)}
{\scriptsize
\setlength{\tabcolsep}{3.8pt}
\renewcommand{\arraystretch}{1.15}
\rowcolors{2}{rowalt}{white}
\begin{xltabular}{\textwidth}{L{1.35cm} L{4.55cm} C{2.05cm} Y}
\toprule
\rowcolor{uibblue!12}\textbf{Lines} & \textbf{Symbol (Class / Function)} & \textbf{Called By} & \textbf{Primary Scientific / Engineering Role} \\
\midrule
\endfirsthead
\toprule
\rowcolor{uibblue!12}\textbf{Lines} & \textbf{Symbol (Class / Function)} & \textbf{Called By} & \textbf{Primary Scientific / Engineering Role} \\
\midrule
\endhead
\bottomrule
\endfoot
L35--L41   & \texttt{get\_best\_device()} & NB 01--05 & Runtime accelerator priority: \texttt{CUDA} $\to$ Apple \texttt{MPS} $\to$ \texttt{CPU}. \\
L44--L49   & \texttt{sync\_device(dev)} & internal (2) & Synchronizes async GPU streams before \texttt{time.perf\_counter()}. \\
L52--L69   & \texttt{resolve\_kiwi\_file(rel\_path)} & internal (3) & Multi-depth local lookup with GitHub \texttt{main} Colab fallback. \\
L72--L86   & \texttt{resolve\_report\_dirs()} & NB 01--05 & Resolves/creates \texttt{figures/} and \texttt{results/} report folders. \\
L94--L124  & \texttt{KiwiDataset} \textit{(dataclass)} & internal (9) & Immutable bundle of 7\,T volume, Micro-GT splits, organ geometry, predicates \& affine. \\
L127--L168 & \texttt{build\_whole\_organ\_geometry()} & NB 05, int. & Berry Otsu mask, $\tilde{r}$, $\tilde{z}$, $R_{\text{ring}}$; supports \texttt{center\_shift\_vox} \& \texttt{radial\_scale}. \\
L171--L231 & \texttt{load\_kiwi\_dataset()} & NB 01--05 & Loads MRI + Micro-GT (\texttt{gt\_nii\_override}), $B_1$ closing envelope, $r_{T_2}$, $B_{\text{non\_dark}}$. \\
L234--L241 & \texttt{save\_aligned\_label\_nifti()} & NB 04 & Exports \texttt{uint16} 3D instance mask bit-aligned to MRI \texttt{sform}/\texttt{qform}. \\
L244--L250 & \texttt{make\_3ch\_tensor()} & internal (2) & Stacks $[\text{vol\_norm}, \text{env\_local}, \text{env\_smooth}]$ into $(1, 3, X, Y, Z)$ tensor. \\
L253--L266 & \texttt{make\_3ch\_and\_gps()} & NB 02, int. & Builds $(\mathbf{X}_3, \mathbf{g})$ with $\mathbf{g} = [\tilde{r}, (\tilde{r}-0.43)^2, \tilde{p}_z]$ and override hooks. \\
L269--L279 & \texttt{build\_distance\_targets()} & internal (1) & Binary foreground + per-instance padded anisotropic 3D EDT in $[0, 1]$. \\
L286--L300 & \texttt{ConvBlock3D} \textit{(nn.Module)} & internal (5) & Standard medical CNN block: $(\text{Conv3d}(3) \to \text{InstanceNorm3d} \to \text{SiLU}) \times 2$. \\
L303--L322 & \texttt{Backbone3DUNet} \textit{(nn.Module)} & internal (6) & 3-level 3D U-Net encoder--decoder ($16 \to 32 \to 64 \to 32 \to 16$ channels). \\
L325--L335 & \texttt{MultiTaskUNet3D} \textit{(nn.Module)} & NB 02, int. & Dual-head 3D U-Net for \texttt{baseline} (\texttt{in\_ch=3}) or naive \texttt{CoordConv3D} (\texttt{in\_ch=6}). \\
L338--L361 & \texttt{AdditiveBayesOrganNet3D} & NB 02, int. & Explicit Bayes log-odds network $z_{\text{app}} + z_{\text{GPS}}$ with auditable \texttt{gps\_log\_prior()}. \\
L364--L388 & \texttt{CoupledOrganNet3D} & NB 02, int. & Feature-concatenation Organ-GPS network ($+864$ params, nonlinear fusion). \\
L395--L410 & \texttt{implication\_violation()} & NB 02 & Fuzzy violation $1 - I(a \Rightarrow b)$ for \texttt{product}, \texttt{lukasiewicz}, or \texttt{godel}. \\
L413--L415 & \texttt{ltn\_universal\_pmean()} & NB 02, int. & LTN generalized $p$-mean quantifier $(\text{mean}(v^p) + 10^{-8})^{1/p}$. \\
L418--L421 & \texttt{soft\_dice\_loss()} & internal (1) & Differentiable smooth Soft Dice loss over foreground logits. \\
L424--L429 & \texttt{supervised\_multitask\_loss()} & internal (3) & Combined $0.45\,\text{BCE}_{w_+=6} + 0.55\,\text{SoftDice} + 2.5\,\text{WeightedMSE}$. \\
L436--L483 & \texttt{pretrain\_ssl\_backbone()} & NB 03--05 & 3D masked-context SSL warm-start on 36 test-disjoint unlabeled cubes. \\
L486--L510 & \texttt{sample\_unlabeled\_blocks()} & NB 03--05 & Stratified 10-latitude unlabeled bank ($\Omega_U$, strictly disjoint from test wedge). \\
L513--L623 & \texttt{train\_neurosymbolic\_variant()} & NB 03--05 & Seed-deterministic training harness for all 4 modes + LTN + columella toggle. \\
L630--L633 & \texttt{build\_ellip()} & internal (1) & Physical-mm 3D ellipsoidal structuring element compensating $\Delta z = 2.13\times\Delta xy$. \\
L636--L657 & \texttt{decode\_deep\_watershed\_3d()} & NB 03, 05, int. & Anisotropic distance-peak 3D watershed + island rescue + min-volume filter. \\
L660--L689 & \texttt{symbolic\_hardcore\_clump\_split()} & NB 03, 04, int. & Declarative Mat\'ern hard-core $K$-means split ($V \ge 2.55\,\text{mm}^3$, $d_{\min} \ge 0.88\,\text{mm}$). \\
L692--L730 & \texttt{evaluate\_semantic\_and\_instance\_3d()} & NB 03--05 & Vectorized 3D IoU matrix, Hungarian matching, $\text{PQ} = \text{SQ} \times \text{RQ}$, Merges/Splits. \\
L733--L750 & \texttt{match\_instance\_sets()} & NB 04 & Returns exact $(\text{TP}_{\text{pred}}, \text{TP}_{\text{gt}}, \text{FN}_{\text{gt}}, \text{FP}_{\text{pred}})$ ID sets for error anatomy. \\
L753--L758 & \texttt{forward\_patch()} & internal (3) & Mode-aware forward dispatcher unifying 1-input and 2-input architectures. \\
L761--L783 & \texttt{calibrate\_threshold\_on\_train()} & NB 03--05 & Leakage-free threshold sweep $\tau \in [0.22, 0.50]$ strictly on labeled training wedge. \\
L786--L847 & \texttt{run\_test\_and\_whole\_fruit()} & NB 03--05 & Unmasked test-wedge evaluation + optional $3{\times}5{\times}8$ tiled 128-slice whole-fruit inference. \\
\end{xltabular}
}

\setcounter{section}{-1}
\section[\texorpdfstring{Module Header, Scientific Stack \& GitHub Endpoint (L1--L29)}{Module Header, Scientific Stack & GitHub Endpoint (L1--L29)}]{Module Header, Scientific Stack \& GitHub Endpoint (L1--L29)}
\begin{annotbox}[Design Notes --- L1--L29: Minimal Pure-PyTorch Dependency Footprint]
\begin{itemize}
    \item \texttt{from \_\_future\_\_ import annotations} (L8) defers PEP~604 union type evaluation (\texttt{np.ndarray | None}) cleanly across Python 3.10--3.12 (both local \texttt{uv} Python 3.12 and Google Colab runtimes).
    \item \textbf{Zero heavy framework lock-in:} The library depends exclusively on standard scientific Python packages (\texttt{torch}, \texttt{nibabel}, \texttt{numpy}, \texttt{scipy}, \texttt{scikit-image}, \texttt{scikit-learn})---no \texttt{MONAI}, \texttt{mmcv}, or custom CUDA kernels---so every notebook runs identically on macOS Apple Silicon (\texttt{mps}), Linux NVIDIA (\texttt{cuda}), and CPU.
    \item \texttt{GITHUB\_RAW\_KIWI} (L28) points explicitly to the \texttt{main} branch so standalone Google Colab sessions automatically download the exact published 7\,T NIfTI volume and Micro-GT files.
\end{itemize}
\end{annotbox}
\begin{srcbox}[{kiwi\_ns\_lib.py $\cdot$ Lines 1--29 --- Module Docstring, Imports \& GitHub Raw Base URL}]
\begin{lstlisting}[style=pyannot,firstnumber=1]
"""kiwi_ns_lib.py — Shared PyTorch & Quantitative Imaging Library for the BMED365
Differentiable Neurosymbolic AI Technical Report ('Thesis + Code' Package).

All 5 companion notebooks in `docs/neurosymbolic-report/notebooks/` import this
module (or download it automatically from GitHub `main` when run in Google Colab),
ensuring zero code drift between experiments, figures, tables, and LaTeX prose.
"""
from __future__ import annotations

import json
import math
import time
import urllib.request
from dataclasses import dataclass
from pathlib import Path

import nibabel as nib
import numpy as np
import torch
import torch.nn as nn
import torch.nn.functional as F
from scipy import ndimage as ndi
from scipy.optimize import linear_sum_assignment
from scipy.spatial.distance import cdist
from skimage import filters, measure, segmentation
from sklearn.cluster import KMeans

GITHUB_RAW_KIWI = "https://raw.githubusercontent.com/arvidl/BMED365-2027/main/Lab4-Comp-Imag/Kiwi-challenge"
\end{lstlisting}
\end{srcbox}

\section[\texorpdfstring{Cross-Platform Runtime \& Portable Path Resolution (L30--L87)}{Cross-Platform Runtime & Portable Path Resolution (L30--L87)}]{Cross-Platform Runtime \& Portable Path Resolution (L30--L87)}
\subsection*{1.1 Accelerator Priority \& Asynchronous GPU Timing (L30--L50)}
\addcontentsline{toc}{subsection}{\texorpdfstring{1.1 Accelerator Priority \& Asynchronous GPU Timing (L30--L50)}{1.1 Accelerator Priority & Asynchronous GPU Timing (L30--L50)}}
\begin{annotbox}[Runtime Notes --- L35--L49: Dynamic Backend Selection \& Wall-Clock Synchronization]
\begin{itemize}
    \item \texttt{get\_best\_device()} (L35--L41) enforces the BMED365 portability rule: never hard-code \texttt{"cuda"} or \texttt{"mps"}. Priority order is \texttt{CUDA} (Colab T4 / Linux workstation) $\to$ Apple \texttt{MPS} (M1--M5 Mac) $\to$ \texttt{CPU}.
    \item \texttt{sync\_device()} (L44--L49) is essential for honest timing in \texttt{pretrain\_ssl\_backbone()} and \texttt{train\_neurosymbolic\_variant()}: both CUDA and Apple Metal execute kernels asynchronously, so calling \texttt{time.perf\_counter()} without draining the command queue would severely under-report GPU training times.
\end{itemize}
\end{annotbox}
\begin{srcbox}[{kiwi\_ns\_lib.py $\cdot$ Lines 30--50 --- Device Dispatch \& GPU Stream Synchronization}]
\begin{lstlisting}[style=pyannot,firstnumber=31]
# =============================================================================
# 1. Cross-Platform Runtime & Portable Path Resolution
# =============================================================================

def get_best_device() -> torch.device:
    """Select fastest available accelerator at runtime: CUDA -> Apple MPS -> CPU."""
    if torch.cuda.is_available():
        return torch.device("cuda")
    if hasattr(torch.backends, "mps") and torch.backends.mps.is_available():
        return torch.device("mps")
    return torch.device("cpu")


def sync_device(dev: torch.device) -> None:
    """Synchronize asynchronous GPU streams before wall-clock timing."""
    if dev.type == "cuda":
        torch.cuda.synchronize()
    elif dev.type == "mps":
        torch.mps.synchronize()
\end{lstlisting}
\end{srcbox}

\subsection*{1.2 Zero-Config Local \& Google Colab File Resolution (L51--L87)}
\addcontentsline{toc}{subsection}{\texorpdfstring{1.2 Zero-Config Local \& Google Colab File Resolution (L51--L87)}{1.2 Zero-Config Local & Google Colab File Resolution (L51--L87)}}
\begin{annotbox}[Engineering Notes --- L52--L86: Working-Directory Independence]
\begin{itemize}
    \item \texttt{resolve\_kiwi\_file()} (L52--L69) probes five relative ancestor depths so notebooks execute cleanly whether JupyterLab is launched from the repository root (\texttt{BMED365-2027/}), from \texttt{Kiwi-challenge/}, or directly inside \texttt{docs/neurosymbolic-report/notebooks/}. If none match (e.g., an ephemeral \texttt{/content} Colab container), it downloads the file once from \texttt{GITHUB\_RAW\_KIWI}.
    \item \texttt{resolve\_report\_dirs()} (L72--L86) anchors all figure PNGs (\texttt{FIG\_DIR}) and JSON/NIfTI outputs (\texttt{RES\_DIR}) into \texttt{docs/neurosymbolic-report/\{figures,results\}/} locally, or local \texttt{./figures} and \texttt{./results} in Colab.
\end{itemize}
\end{annotbox}
\begin{srcbox}[{kiwi\_ns\_lib.py $\cdot$ Lines 51--87 --- Portable Data Lookup \& Report Output Directory Resolution}]
\begin{lstlisting}[style=pyannot,firstnumber=52]
def resolve_kiwi_file(rel_path: str) -> Path:
    """Resolve a Kiwi-challenge data/asset file locally or auto-download on Colab."""
    here = Path.cwd()
    candidates = [
        here / rel_path,
        here / ".." / rel_path,
        here / ".." / ".." / rel_path,
        here / ".." / ".." / ".." / rel_path,
        here / "Lab4-Comp-Imag/Kiwi-challenge" / rel_path,
    ]
    for cand in candidates:
        if cand.exists():
            return cand.resolve()
    dst = Path(rel_path)
    dst.parent.mkdir(parents=True, exist_ok=True)
    print(f"Downloading {rel_path} from GitHub main for Colab session...")
    urllib.request.urlretrieve(f"{GITHUB_RAW_KIWI}/{rel_path}", dst)
    return dst.resolve()


def resolve_report_dirs() -> tuple[Path, Path]:
    """Return (figures_dir, results_dir) for the neurosymbolic-report package."""
    here = Path.cwd()
    for base in [here / "..", here, here / "Lab4-Comp-Imag/Kiwi-challenge/docs/neurosymbolic-report"]:
        if (base / "notebooks").exists() or base.name == "neurosymbolic-report":
            figs = (base / "figures").resolve()
            res = (base / "results").resolve()
            figs.mkdir(parents=True, exist_ok=True)
            res.mkdir(parents=True, exist_ok=True)
            return figs, res
    figs = (here / "figures").resolve()
    res = (here / "results").resolve()
    figs.mkdir(parents=True, exist_ok=True)
    res.mkdir(parents=True, exist_ok=True)
    return figs, res
\end{lstlisting}
\end{srcbox}

\section[\texorpdfstring{Data Loading, Whole-Organ Coordinate Frame \& Biophysical Predicates (L88--L280)}{Data Loading, Whole-Organ Coordinate Frame & Biophysical Predicates (L88--L280)}]{Data Loading, Whole-Organ Coordinate Frame \& Biophysical Predicates (L88--L280)}
\subsection*{2.1 The \texttt{KiwiDataset} Container (L88--L125)}
\addcontentsline{toc}{subsection}{\texorpdfstring{2.1 The \texttt{KiwiDataset} Container (L88--L125)}{2.1 The KiwiDataset Container (L88--L125)}}
\begin{annotbox}[Data Contract --- L93--L124: Unified 3D Volume, Split \& Predicate Bundle]
\texttt{KiwiDataset} encapsulates all 27 arrays and scalar metadata produced by \texttt{load\_kiwi\_dataset()} so downstream functions never pass loose tuples of 15 arrays:
\begin{itemize}
    \item \textbf{Spatial ROI \& Leakage-Free Split:} \texttt{roi\_gt} is the $80 \times 80 \times 16$ equatorial Micro-GT sector (\texttt{rx0:rx1}, \texttt{ry0:ry1}, \texttt{rz0:rz1} $= [128,208)\times[128,208)\times[56,72)$); \texttt{split\_y = 40} partitions it into voxel-disjoint \texttt{train\_gt} ($y \in [128, 168)$, 59 seeds) and \texttt{test\_gt} ($y \in [168, 208)$, 52 seeds).
    \item \textbf{Central Columella Hard-Negative Box (L119--L122):} \texttt{cx0..cx1 = 112..144}, \texttt{cy0..cy1 = 112..144} carves a $32 \times 32 \times 16$ seed-free central-core crop used in \texttt{train\_neurosymbolic\_variant()} and toggled causally in Notebook~04.
    \item \textbf{Affine Provenance (L123--L124):} Stores the original $4 \times 4$ NIfTI \texttt{affine} and \texttt{gt\_affine\_max\_diff} ($\max|\mathbf{A}_{\text{MRI}} - \mathbf{A}_{\text{GT}}| = 0.0$) for bit-aligned Freeview/ITK-SNAP verification.
\end{itemize}
\end{annotbox}
\begin{srcbox}[{kiwi\_ns\_lib.py $\cdot$ Lines 88--125 --- KiwiDataset Dataclass Specification}]
\begin{lstlisting}[style=pyannot,firstnumber=89]
# =============================================================================
# 2. Data Loading, Whole-Organ Coordinate Frame & Biophysical Predicates
# =============================================================================

@dataclass
class KiwiDataset:
    vol: np.ndarray
    spacing: tuple[float, float, float]
    vox_vol_mm3: float
    gt_full: np.ndarray
    roi_gt: np.ndarray
    train_gt: np.ndarray
    test_gt: np.ndarray
    rx0: int
    rx1: int
    ry0: int
    ry1: int
    rz0: int
    rz1: int
    split_y: int
    fruit_mask: np.ndarray
    r_norm: np.ndarray
    z_norm: np.ndarray
    pole_dist: np.ndarray
    soft_ring: np.ndarray
    vol_norm: np.ndarray
    env_local: np.ndarray
    env_smooth: np.ndarray
    t2_ratio: np.ndarray
    non_dark: np.ndarray
    cx0: int = 112
    cx1: int = 144
    cy0: int = 112
    cy1: int = 144
    affine: np.ndarray | None = None
    gt_affine_max_diff: float = 0.0
\end{lstlisting}
\end{srcbox}

\subsection*{2.2 Intrinsic Organ Coordinate Frame \& Stress-Test Hooks (L126--L169)}
\addcontentsline{toc}{subsection}{\texorpdfstring{2.2 Intrinsic Organ Coordinate Frame \& Stress-Test Hooks (L126--L169)}{2.2 Intrinsic Organ Coordinate Frame & Stress-Test Hooks (L126--L169)}}
\begin{annotbox}[Mathematical Annotation --- L127--L168: Slice-Normalized Radius \& Smooth Compartment Predicate]
\texttt{build\_whole\_organ\_geometry()} maps the raw 7\,T volume into an atlas-free intrinsic coordinate frame invariant to the fruit's polar taper:
\begin{enumerate}
    \item \textbf{Intact Berry Mask (L139--L146):} Anisotropic Gaussian smoothing ($\sigma = 1.2\,\text{mm}$ converted to voxels via $1.2 / s_k$), Otsu thresholding at $0.35\times$, 2D slice-wise + 3D hole filling, and largest connected component isolate \texttt{fruit\_mask} ($V_{\text{kiwi}} = 85.97\,\text{cm}^3$).
    \item \textbf{Area-Equivalent Radial Depth (L149--L156):} For every slice $z$ with $\ge 80$ berry voxels, $R_{\text{fruit}}(z) = \sqrt{\text{Area}(z)/\pi} \cdot \alpha$ and $\tilde{r}(x,y,z) = \sqrt{(x - (c_x + \Delta x))^2 + (y - (c_y + \Delta y))^2} / R_{\text{fruit}}(z)$. Outside-fruit voxels default to $1.5$.
    \item \textbf{Smooth Differentiable Inner-Pericarp Membership $R_{\text{ring}} \in [0, 1]$ (L162--L167):}
    \[
        R_{\text{ring}}(\mathbf{x}) = \sigma\!\left(\frac{\tilde{r} - 0.22}{0.02}\right)\sigma\!\left(\frac{0.65 - \tilde{r}}{0.02}\right) \cdot \operatorname{clip}\!\left(\frac{z - 16}{5}, 0, 1\right)\operatorname{clip}\!\left(\frac{111 - z}{5}, 0, 1\right) \cdot M_{\text{fruit}}(\mathbf{x}).
    \]
    \item \textbf{Pathological Stress-Test Parameters:} Optional \texttt{center\_shift\_vox} $(\Delta x, \Delta y, \Delta z)$ and \texttt{radial\_scale} $\alpha$ let Notebook~05 inject synthetic clinical mass effect ($0$--$8\,\text{mm}$) and atrophy/hypertrophy ($\alpha \in [0.80, 1.25]$) without modifying any core code.
\end{enumerate}
\end{annotbox}
\begin{srcbox}[{kiwi\_ns\_lib.py $\cdot$ Lines 126--169 --- build\_whole\_organ\_geometry()}]
\begin{lstlisting}[style=pyannot,firstnumber=127]
def build_whole_organ_geometry(
    vol_3d: np.ndarray,
    sp: tuple[float, float, float],
    center_shift_vox: tuple[float, float, float] = (0.0, 0.0, 0.0),
    radial_scale: float = 1.0,
) -> tuple[np.ndarray, np.ndarray, np.ndarray, np.ndarray]:
    """Extract intact berry mask, normalized radial coordinate r/R_fruit(z),
    axial coordinate z/Z, and smooth differentiable inner-pericarp membership R_ring.

    Optional `center_shift_vox` (dx, dy, dz) and `radial_scale` allow synthetic
    pathological distortion / mass-effect stress testing (Notebook 05).
    """
    v_sm = ndi.gaussian_filter(vol_3d, sigma=(1.2 / sp[0], 1.2 / sp[1], 1.2 / sp[2]))
    fg_raw = v_sm > (filters.threshold_otsu(v_sm[::2, ::2, ::2]) * 0.35)
    fg_fill = np.zeros_like(fg_raw)
    for z in range(vol_3d.shape[2]):
        fg_fill[:, :, z] = ndi.binary_fill_holes(fg_raw[:, :, z])
    fg_fill = ndi.binary_fill_holes(fg_fill)
    lbl_fg, _ = ndi.label(fg_fill)
    fruit_mask = lbl_fg == (1 + int(np.argmax(np.bincount(lbl_fg.ravel())[1:])))

    dx, dy, dz = center_shift_vox
    r_norm_3d = np.full(vol_3d.shape, 1.5, dtype=np.float32)
    yy, xx = np.meshgrid(np.arange(vol_3d.shape[1]), np.arange(vol_3d.shape[0]))
    for z in range(vol_3d.shape[2]):
        m = fruit_mask[:, :, z]
        if m.sum() >= 80:
            cx, cy = float(xx[m].mean()) + dx, float(yy[m].mean()) + dy
            r_fruit = max(np.sqrt(m.sum() / np.pi) * radial_scale, 1.0)
            r_norm_3d[:, :, z] = np.sqrt((xx - cx) ** 2 + (yy - cy) ** 2) / r_fruit

    z_idx = np.arange(vol_3d.shape[2], dtype=np.float32) + dz
    z_norm_1d = np.clip(z_idx / float(vol_3d.shape[2] - 1), 0.0, 1.0)
    z_norm_3d = np.broadcast_to(z_norm_1d[None, None, :], vol_3d.shape).copy()

    z_gate = np.clip((z_idx[None, None, :] - 16.0) / 5.0, 0.0, 1.0) * np.clip(
        (111.0 - z_idx[None, None, :]) / 5.0, 0.0, 1.0
    )
    r_in = 1.0 / (1.0 + np.exp(-np.clip((r_norm_3d - 0.22) / 0.02, -20.0, 20.0)))
    r_out = 1.0 / (1.0 + np.exp(-np.clip((0.65 - r_norm_3d) / 0.02, -20.0, 20.0)))
    soft_ring = (r_in * r_out * z_gate * fruit_mask.astype(np.float32)).astype(np.float32)
    return fruit_mask, r_norm_3d, z_norm_3d, soft_ring
\end{lstlisting}
\end{srcbox}

\subsection*{2.3 7\,T MRI Loading, $B_1$-Invariant Morphological Envelope \& Bit-Aligned NIfTI Export (L170--L242)}
\addcontentsline{toc}{subsection}{\texorpdfstring{2.3 7\,T MRI Loading, $B_1$-Invariant Morphological Envelope \& Bit-Aligned NIfTI Export (L170--L242)}{2.3 7 T MRI Loading, B1-Invariant Morphological Envelope & Bit-Aligned NIfTI Export (L170--L242)}}
\begin{annotbox}[Physics \& Provenance Notes --- L171--L241: $B_1$ Shading Cancellation \& Manual GT Override]
\begin{itemize}
    \item \textbf{Manual Ground-Truth Override (L171, L179--L183):} Passing \texttt{gt\_nii\_override="MyStuff/kiwi-gt/kiwi\_gt\_manual\_v1.nii.gz"} re-runs all five notebooks against a manually refined \texttt{uint16} mask while keeping the repository's shared silver standard untouched.
    \item \textbf{Safe Integer Rounding (L192):} \texttt{np.rint(np.asarray(gt\_img.dataobj)).astype(np.uint16)} prevents floating-point NIfTI scale-slope artifacts from corrupting integer instance IDs.
    \item \textbf{Removing 7\,T RF $B_1$ Shading via 3D Morphological Closing (L207--L221):} Normalizing by the 99th inside-fruit percentile (\texttt{vol\_norm}), then applying \texttt{ndi.grey\_closing(size=(7, 7, 3))} ($\approx 1.64 \times 1.64 \times 1.50\,\text{mm}^3$) bridges across dark $\sim 1.5\,\text{mm}$ seeds to reconstruct the bright surrounding mucilage gel envelope \texttt{env\_local}. Dividing gives the $B_1$-invariant local contrast ratio $r_{T_2} = \operatorname{clip}(I / (E_{\text{local}} + 0.03), 0, 1.25)$, and the smooth physics violation predicate:
    \[
        B_{\text{non\_dark}}(\mathbf{x}) = \max\!\left(\operatorname{clip}\!\left(\frac{r_{T_2} - 0.82}{0.15},\, 0,\, 1\right),\; \operatorname{clip}\!\left(\frac{0.32 - E_{\text{local}}}{0.10},\, 0,\, 1\right)\right).
    \]
    The $0.82$ onset protects $100\%$ of seed cores and $93.4\%$ of 1-voxel partial-volume rims while penalizing bright gel ($r_{T_2} \to 1$) and dry columella/air ($E_{\text{local}} < 0.32$).
    \item \textbf{Bit-Aligned Export (L234--L241):} \texttt{save\_aligned\_label\_nifti()} sets both \texttt{sform} and \texttt{qform} to \texttt{code=1} with the exact MRI affine so FreeSurfer \texttt{freeview}, 3D Slicer, and ITK-SNAP superimpose predictions with zero resampling.
\end{itemize}
\end{annotbox}
\begin{srcbox}[{kiwi\_ns\_lib.py $\cdot$ Lines 170--242 --- load\_kiwi\_dataset() \& save\_aligned\_label\_nifti()}]
\begin{lstlisting}[style=pyannot,firstnumber=171]
def load_kiwi_dataset(gt_nii_override: str | Path | None = None) -> KiwiDataset:
    """Load 7T RARE kiwi MRI, curated 3D Micro-GT, organ coordinates & biophysical predicates.

    Pass `gt_nii_override` (e.g. 'MyStuff/kiwi-gt/kiwi_gt_manual_v1.nii.gz') to audit,
    evaluate, or retrain against a manually refined 3D instance mask without modifying
    the shared silver standard in `data/kiwi_micro_gt_labels.nii.gz`.
    """
    nii_path = resolve_kiwi_file("data/kiwi_20050217191129_131073_RARE_high_res_axial.nii.gz")
    gt_nii_path = (
        Path(gt_nii_override).expanduser().resolve()
        if gt_nii_override is not None
        else resolve_kiwi_file("data/kiwi_micro_gt_labels.nii.gz")
    )
    gt_json_path = resolve_kiwi_file("data/kiwi_micro_gt_roi.json")

    nii_img = nib.load(str(nii_path))
    vol = nii_img.get_fdata(dtype=np.float32)
    spacing = tuple(float(s) for s in nii_img.header.get_zooms()[:3])
    vox_vol_mm3 = float(np.prod(spacing))

    gt_img = nib.load(str(gt_nii_path))
    gt_full = np.rint(np.asarray(gt_img.dataobj)).astype(np.uint16)
    gt_affine_max_diff = float(np.max(np.abs(nii_img.affine - gt_img.affine)))
    gt_meta = json.loads(gt_json_path.read_text(encoding="utf-8"))
    rx0, rx1 = gt_meta["roi_bounds_voxel_xyz"]["x"]
    ry0, ry1 = gt_meta["roi_bounds_voxel_xyz"]["y"]
    rz0, rz1 = gt_meta["roi_bounds_voxel_xyz"]["z"]
    roi_gt = gt_full[rx0:rx1, ry0:ry1, rz0:rz1]

    split_y = 40
    train_gt = roi_gt[:, :split_y, :].copy()
    test_gt = roi_gt[:, split_y:, :].copy()

    fruit_mask, r_norm, z_norm, soft_ring = build_whole_organ_geometry(vol, spacing)
    pole_dist = (np.abs(z_norm - 0.5) * 2.0).astype(np.float32)

    p99 = float(np.percentile(vol[fruit_mask], 99.0))
    vol_norm = np.clip(vol / (p99 + 1e-6), 0.0, 1.25).astype(np.float32)
    env_close = ndi.grey_closing(vol_norm, size=(7, 7, 3))
    env_local = ndi.gaussian_filter(
        env_close, sigma=(0.8 / spacing[0], 0.8 / spacing[1], 0.8 / spacing[2])
    ).astype(np.float32)
    env_smooth = ndi.gaussian_filter(
        vol_norm, sigma=(2.0 / spacing[0], 2.0 / spacing[1], 1.6 / spacing[2])
    ).astype(np.float32)

    t2_ratio = np.clip(vol_norm / (env_local + 0.03), 0.0, 1.25).astype(np.float32)
    non_dark = np.maximum(
        np.clip((t2_ratio - 0.82) / 0.15, 0.0, 1.0),
        np.clip((0.32 - env_local) / 0.10, 0.0, 1.0),
    ).astype(np.float32)

    return KiwiDataset(
        vol=vol, spacing=spacing, vox_vol_mm3=vox_vol_mm3,
        gt_full=gt_full, roi_gt=roi_gt, train_gt=train_gt, test_gt=test_gt,
        rx0=rx0, rx1=rx1, ry0=ry0, ry1=ry1, rz0=rz0, rz1=rz1, split_y=split_y,
        fruit_mask=fruit_mask, r_norm=r_norm, z_norm=z_norm, pole_dist=pole_dist,
        soft_ring=soft_ring, vol_norm=vol_norm, env_local=env_local,
        env_smooth=env_smooth, t2_ratio=t2_ratio, non_dark=non_dark,
        affine=nii_img.affine.copy(), gt_affine_max_diff=gt_affine_max_diff,
    )


def save_aligned_label_nifti(lbl_3d: np.ndarray, ref_affine: np.ndarray, out_path: Path) -> Path:
    """Save a uint16 3D instance label mask bit-aligned to the original 7T MRI affine."""
    out_path.parent.mkdir(parents=True, exist_ok=True)
    img = nib.Nifti1Image(lbl_3d.astype(np.uint16), ref_affine)
    img.set_sform(ref_affine, code=1)
    img.set_qform(ref_affine, code=1)
    nib.save(img, str(out_path))
    return out_path
\end{lstlisting}
\end{srcbox}

\subsection*{2.4 Appearance/GPS Tensor Builders \& Per-Instance Anisotropic 3D EDT (L243--L280)}
\addcontentsline{toc}{subsection}{\texorpdfstring{2.4 Appearance/GPS Tensor Builders \& Per-Instance Anisotropic 3D EDT (L243--L280)}{2.4 Appearance/GPS Tensor Builders & Per-Instance Anisotropic 3D EDT (L243--L280)}}
\begin{annotbox}[Tensor Engineering Notes --- L244--L279: Quadratic Radial Bowl \& Padded EDT]
\begin{itemize}
    \item \texttt{make\_3ch\_tensor()} (L244--L250) packages three complementary appearance channels $[I,\, E_{\text{local}},\, E_{\text{smooth}}]$ so the CNN sees raw $T_2$ intensity, local mucilage context, and coarse $B_1$ shading simultaneously.
    \item \texttt{make\_3ch\_and\_gps()} (L253--L266) builds the 3-channel un-normalized \texttt{Organ-GPS} tensor $[\tilde{r},\; (\tilde{r} - 0.43)^2,\; |2z/Z - 1|]$. Providing the explicit quadratic bowl $(\tilde{r} - 0.43)^2$ centered at mid-pericarp lets a tiny 2-layer MLP represent a smooth annular ring cleanly without needing high capacity. Optional \texttt{r\_override}/\texttt{pole\_override} enable test-time coordinate perturbation in Notebook~05.
    \item \texttt{build\_distance\_targets()} (L269--L279) computes each seed's normalized center-distance dome $y_d \in [0, 1]$ via \texttt{ndi.distance\_transform\_edt(..., sampling=sp)} inside a \textbf{1-voxel zero-padded bounding box} (\texttt{pad\_width=1}). Padding is critical: without it, seeds touching the ROI boundary would have artificially nonzero distance at the crop edge, corrupting watershed peaks.
\end{itemize}
\end{annotbox}
\begin{srcbox}[{kiwi\_ns\_lib.py $\cdot$ Lines 243--280 --- Tensor Packaging \& Per-Instance Anisotropic 3D EDT}]
\begin{lstlisting}[style=pyannot,firstnumber=244]
def make_3ch_tensor(ds: KiwiDataset, x0: int, x1: int, y0: int, y1: int, z0: int, z1: int) -> torch.Tensor:
    """Stack [vol_norm, env_local, env_smooth] -> (1, 3, X, Y, Z)."""
    arr = np.stack(
        [ds.vol_norm[x0:x1, y0:y1, z0:z1], ds.env_local[x0:x1, y0:y1, z0:z1], ds.env_smooth[x0:x1, y0:y1, z0:z1]],
        axis=0,
    )[None, ...]
    return torch.from_numpy(arr)


def make_3ch_and_gps(
    ds: KiwiDataset,
    x0: int, x1: int, y0: int, y1: int, z0: int, z1: int,
    r_override: np.ndarray | None = None,
    pole_override: np.ndarray | None = None,
) -> tuple[torch.Tensor, torch.Tensor]:
    """Return (x3, gps) where gps = [r/R, (r/R - 0.43)^2, |2z/Z - 1|]."""
    x3 = make_3ch_tensor(ds, x0, x1, y0, y1, z0, z1)
    r_src = ds.r_norm if r_override is None else r_override
    p_src = ds.pole_dist if pole_override is None else pole_override
    rn = np.clip(r_src[x0:x1, y0:y1, z0:z1], 0.0, 1.5)
    zn = p_src[x0:x1, y0:y1, z0:z1]
    gps = np.stack([rn, (rn - 0.43) ** 2, zn], axis=0)[None, ...].astype(np.float32)
    return x3, torch.from_numpy(gps)


def build_distance_targets(inst_lbl: np.ndarray, sp: tuple[float, float, float]) -> tuple[np.ndarray, np.ndarray]:
    """Compute binary semantic foreground + per-instance normalized 3D EDT in [0, 1]."""
    sem_fg = (inst_lbl > 0).astype(np.float32)
    norm_dist = np.zeros(inst_lbl.shape, dtype=np.float32)
    for prop in measure.regionprops(inst_lbl):
        s0, s1, s2 = prop.slice
        sub = np.pad(inst_lbl[s0, s1, s2] == prop.label, pad_width=1, mode="constant", constant_values=False)
        d_local = ndi.distance_transform_edt(sub, sampling=sp)[1:-1, 1:-1, 1:-1]
        m_local = sub[1:-1, 1:-1, 1:-1]
        norm_dist[s0, s1, s2][m_local] = (d_local[m_local] / max(float(d_local.max()), 1e-6)).astype(np.float32)
    return sem_fg, norm_dist
\end{lstlisting}
\end{srcbox}

\section[\texorpdfstring{Neural \& Differentiable Neurosymbolic 3D Architectures (L281--L389)}{Neural & Differentiable Neurosymbolic 3D Architectures (L281--L389)}]{Neural \& Differentiable Neurosymbolic 3D Architectures (L281--L389)}
\subsection*{3.1 Shared 3D U-Net Backbone \& Pure/Naive Multi-Task Variants (L281--L336)}
\addcontentsline{toc}{subsection}{\texorpdfstring{3.1 Shared 3D U-Net Backbone \& Pure/Naive Multi-Task Variants (L281--L336)}{3.1 Shared 3D U-Net Backbone & Pure/Naive Multi-Task Variants (L281--L336)}}
\begin{annotbox}[Architecture Notes --- L286--L335: Why \texttt{InstanceNorm3d} Helps Texture Yet Breaks \texttt{CoordConv3D}]
\begin{itemize}
    \item \texttt{ConvBlock3D} (L286--L300) omits convolution bias (\texttt{bias=False}) because the immediately following \texttt{InstanceNorm3d(out\_c, affine=True)} subtracts the per-channel mean and supplies its own learnable shift $\beta_c$. \texttt{SiLU} provides smooth non-zero gradients everywhere.
    \item \texttt{Backbone3DUNet} (L303--L322) is a symmetric 3-level U-Net ($16 \to 32 \to 64 \to 32 \to 16$, two $2{\times}2{\times}2$ \texttt{MaxPool3d} downsamples and two \texttt{ConvTranspose3d} upsamples, requiring input spatial dims divisible by $4$). Keeping channel width compact (\texttt{base\_ch=16}, $\sim 340\text{k}$ params) prevents memorization on $59$ labeled training seeds and trains in $\sim 15$\,s on Apple MPS or Colab T4.
    \item \texttt{MultiTaskUNet3D} (L325--L335) serves both \textbf{Model~A} (\texttt{in\_ch=3}, pure neural baseline) and \textbf{Model~B} (\texttt{in\_ch=6}, naive \texttt{CoordConv3D})---the latter fails catastrophically because \texttt{enc1}'s \texttt{InstanceNorm3d} erases absolute patch coordinates.
\end{itemize}
\end{annotbox}
\begin{srcbox}[{kiwi\_ns\_lib.py $\cdot$ Lines 281--336 --- ConvBlock3D, Backbone3DUNet \& MultiTaskUNet3D}]
\begin{lstlisting}[style=pyannot,firstnumber=282]
# =============================================================================
# 3. Neural & Differentiable Neurosymbolic 3D Architectures
# =============================================================================

class ConvBlock3D(nn.Module):
    """Standard 3D medical CNN block (Conv3d -> InstanceNorm3d -> SiLU, x2)."""
    def __init__(self, in_c: int, out_c: int):
        super().__init__()
        self.net = nn.Sequential(
            nn.Conv3d(in_c, out_c, kernel_size=3, padding=1, bias=False),
            nn.InstanceNorm3d(out_c, affine=True),
            nn.SiLU(inplace=True),
            nn.Conv3d(out_c, out_c, kernel_size=3, padding=1, bias=False),
            nn.InstanceNorm3d(out_c, affine=True),
            nn.SiLU(inplace=True),
        )

    def forward(self, x: torch.Tensor) -> torch.Tensor:
        return self.net(x)


class Backbone3DUNet(nn.Module):
    """Compact 3-level 3D U-Net (~340k params) shared across all benchmarks."""
    def __init__(self, in_ch: int = 3, base_ch: int = 16):
        super().__init__()
        self.enc1 = ConvBlock3D(in_ch, base_ch)
        self.pool1 = nn.MaxPool3d(2)
        self.enc2 = ConvBlock3D(base_ch, base_ch * 2)
        self.pool2 = nn.MaxPool3d(2)
        self.bot = ConvBlock3D(base_ch * 2, base_ch * 4)
        self.up2 = nn.ConvTranspose3d(base_ch * 4, base_ch * 2, kernel_size=2, stride=2)
        self.dec2 = ConvBlock3D(base_ch * 4, base_ch * 2)
        self.up1 = nn.ConvTranspose3d(base_ch * 2, base_ch, kernel_size=2, stride=2)
        self.dec1 = ConvBlock3D(base_ch * 2, base_ch)

    def forward(self, x: torch.Tensor) -> torch.Tensor:
        e1 = self.enc1(x)
        e2 = self.enc2(self.pool1(e1))
        b = self.bot(self.pool2(e2))
        d2 = self.dec2(torch.cat([self.up2(b), e2], dim=1))
        return self.dec1(torch.cat([self.up1(d2), e1], dim=1))


class MultiTaskUNet3D(nn.Module):
    """Pure neural (in_ch=3) or naive CoordConv3D+InstanceNorm3d (in_ch=6) 3D U-Net."""
    def __init__(self, in_ch: int = 3, base_ch: int = 16):
        super().__init__()
        self.backbone = Backbone3DUNet(in_ch, base_ch)
        self.head_sem = nn.Conv3d(base_ch, 1, kernel_size=1)
        self.head_dist = nn.Conv3d(base_ch, 1, kernel_size=1)

    def forward(self, x: torch.Tensor):
        feat = self.backbone(x)
        return self.head_sem(feat), torch.sigmoid(self.head_dist(feat))
\end{lstlisting}
\end{srcbox}

\subsection*{3.2 Un-Normalized Bayes Log-Odds \& Coupled \texttt{Organ-GPS} Networks (L337--L389)}
\addcontentsline{toc}{subsection}{\texorpdfstring{3.2 Un-Normalized Bayes Log-Odds \& Coupled \texttt{Organ-GPS} Networks (L337--L389)}{3.2 Un-Normalized Bayes Log-Odds & Coupled Organ-GPS Networks (L337--L389)}}
\begin{annotbox}[Mathematical Annotation --- L338--L388: Additive Bayes Log-Odds vs.\ Coupled Feature Fusion]
\begin{itemize}
    \item \textbf{\texttt{AdditiveBayesOrganNet3D} (L338--L361, $+353$ params):} Isolates $\mathbf{g}$ in a purely point-wise ($1{\times}1{\times}1$) MLP ($3 \to 16 \to 16 \to 1$) with \textbf{zero spatial normalization}:
    \[
        z_{\text{total}}(\mathbf{x}) = z_{\text{app}}\bigl(\mathbf{f}_\theta(\mathbf{X}_3)\bigr) + z_{\text{GPS}}\bigl(\mathbf{g}(\mathbf{x}); \psi\bigr), \qquad
        (3\cdot 16 + 16) + (16\cdot 16 + 16) + (16\cdot 1 + 1) = 64 + 272 + 17 = 353\text{ params.}
    \]
    Because \texttt{gps\_log\_prior()} (L356--L357) is callable standalone on any synthetic $(\tilde{r}, \tilde{p}_z)$ mesh, Notebook~04 audits the exact learned spatial odds veto without approximation. Under mass effect (Notebook~05), finite $z_{\text{GPS}}$ lets strong MRI appearance $z_{\text{app}}$ override stale coordinates smoothly.
    \item \textbf{\texttt{CoupledOrganNet3D} (L364--L388, $+864$ params):} Concatenates the 16-channel un-normalized GPS embedding $\mathbf{g}$ with the 16-channel decoder feature map $\mathbf{f}$ before a 2-layer $1{\times}1{\times}1$ semantic head ($32 \to 16 \to 1$), permitting nonlinear texture--location interactions at the cost of losing exact additive decomposability.
    \item Notice that in both models, \texttt{head\_dist} depends strictly on local appearance features $\mathbf{f}$ (distance-to-seed-center is purely local geometry).
\end{itemize}
\end{annotbox}
\begin{srcbox}[{kiwi\_ns\_lib.py $\cdot$ Lines 337--389 --- AdditiveBayesOrganNet3D \& CoupledOrganNet3D}]
\begin{lstlisting}[style=pyannot,firstnumber=338]
class AdditiveBayesOrganNet3D(nn.Module):
    """Explicit Bayes Log-Odds Neurosymbolic 3D Network:
      z_total(x) = z_app(f_theta(X)) + z_gps(h_psi(gps)).
    Directly auditable spatial log-prior (+353 parameters over MultiTaskUNet3D).
    """
    def __init__(self, in_ch: int = 3, gps_ch: int = 3, base_ch: int = 16):
        super().__init__()
        self.backbone = Backbone3DUNet(in_ch, base_ch)
        self.head_app = nn.Conv3d(base_ch, 1, kernel_size=1)
        self.gps_mlp = nn.Sequential(
            nn.Conv3d(gps_ch, 16, kernel_size=1),
            nn.SiLU(),
            nn.Conv3d(16, 16, kernel_size=1),
            nn.SiLU(),
        )
        self.head_gps = nn.Conv3d(16, 1, kernel_size=1)
        self.head_dist = nn.Conv3d(base_ch, 1, kernel_size=1)

    def gps_log_prior(self, gps: torch.Tensor) -> torch.Tensor:
        return self.head_gps(self.gps_mlp(gps))

    def forward(self, x3: torch.Tensor, gps: torch.Tensor):
        f = self.backbone(x3)
        return self.head_app(f) + self.gps_log_prior(gps), torch.sigmoid(self.head_dist(f))


class CoupledOrganNet3D(nn.Module):
    """Coupled feature-fusion Organ-GPS 3D Network (NB 06 architecture):
      z_total(x) = head_sem([f_theta(X), h_psi(gps)]).
    Allows nonlinear interaction between appearance and organ position (+864 params).
    """
    def __init__(self, in_ch: int = 3, gps_ch: int = 3, base_ch: int = 16):
        super().__init__()
        self.backbone = Backbone3DUNet(in_ch, base_ch)
        self.gps_mlp = nn.Sequential(
            nn.Conv3d(gps_ch, 16, kernel_size=1),
            nn.SiLU(),
            nn.Conv3d(16, 16, kernel_size=1),
            nn.SiLU(),
        )
        self.head_sem = nn.Sequential(
            nn.Conv3d(base_ch + 16, 16, kernel_size=1),
            nn.SiLU(),
            nn.Conv3d(16, 1, kernel_size=1),
        )
        self.head_dist = nn.Conv3d(base_ch, 1, kernel_size=1)

    def forward(self, x3: torch.Tensor, gps: torch.Tensor):
        f = self.backbone(x3)
        g = self.gps_mlp(gps)
        return self.head_sem(torch.cat([f, g], dim=1)), torch.sigmoid(self.head_dist(f))
\end{lstlisting}
\end{srcbox}

\section[\texorpdfstring{Differentiable Real Logic (LTN) \& Supervised Loss Primitives (L390--L430)}{Differentiable Real Logic (LTN) & Supervised Loss Primitives (L390--L430)}]{Differentiable Real Logic (LTN) \& Supervised Loss Primitives (L390--L430)}
\begin{annotbox}[Mathematical Annotation --- L395--L429: Smooth Fuzzy Implications, $p$-Mean Quantifiers \& Multi-Task Balance]
\begin{itemize}
    \item \texttt{implication\_violation()} (L395--L410) computes $v(a, b) = 1 - I(a \Rightarrow b) \in [0, 1]$ for three fuzzy logics:
    \begin{itemize}
        \item \texttt{"product"} (Reichenbach): $v = a(1 - b)$, smooth everywhere with zero dead zones (used in all training).
        \item \texttt{"lukasiewicz"}: $v = \max(0, a - b)$, piecewise linear with a dead triangle where $a \le b$.
        \item \texttt{"godel"}: $v = (1 - b)\,\mathbf{1}\{a > b\}$, discontinuous along $a = b$.
    \end{itemize}
    \item \texttt{ltn\_universal\_pmean()} (L413--L415) implements $\mathcal{L}_p(\mathbf{v}) = (\frac{1}{N}\sum_i v_i^p + 10^{-8})^{1/p}$. Adding $\varepsilon = 10^{-8}$ \emph{inside} the $1/p$-th root prevents \texttt{NaN} gradients ($\frac{d}{dx} x^{1/p} \to \infty$ as $x \to 0^+$) whenever a batch achieves near-zero violation. At $p = 3$, rare $0.8\%$ polar FPs receive a $24.98\times$ stronger gradient than plain averaging ($p = 1$).
    \item \texttt{supervised\_multitask\_loss()} (L424--L429) balances severe voxel imbalance (seeds occupy $<10\%$ even inside the ROI) using \texttt{pos\_weight=6.0} in BCE, smooth \texttt{soft\_dice\_loss}, and foreground-weighted distance MSE:
    \[
        \mathcal{L}_{\text{sup}} = 0.45\,\text{BCE}_{w_+=6} + 0.55\,\mathcal{L}_{\text{SoftDice}} + 2.5\,\mathbb{E}\bigl[(1 + 6 y_s)(\hat{d} - y_d)^2\bigr].
    \]
\end{itemize}
\end{annotbox}
\begin{srcbox}[{kiwi\_ns\_lib.py $\cdot$ Lines 390--430 --- Fuzzy Implication, LTN p-Mean \& Supervised Multi-Task Losses}]
\begin{lstlisting}[style=pyannot,firstnumber=391]
# =============================================================================
# 4. Differentiable Real Logic (LTN) & Supervised Loss Primitives
# =============================================================================

def implication_violation(premise: torch.Tensor, conclusion: torch.Tensor, family: str = "product") -> torch.Tensor:
    """Compute fuzzy implication violation v = 1 - I(premise -> conclusion) in [0, 1].
    Families:
      - 'product' (Reichenbach): v = premise * (1 - conclusion)   [smooth everywhere]
      - 'lukasiewicz'          : v = max(0, premise - conclusion) [piecewise linear, dead zone]
      - 'godel'                : v = (1 - conclusion) if premise > conclusion else 0 [discontinuous]
    """
    a = premise.clamp(0.0, 1.0)
    b = conclusion.clamp(0.0, 1.0)
    if family == "product":
        return a * (1.0 - b)
    if family == "lukasiewicz":
        return torch.clamp(a - b, min=0.0, max=1.0)
    if family == "godel":
        return torch.where(a > b, 1.0 - b, torch.zeros_like(a))
    raise ValueError(f"Unknown fuzzy implication family: {family}")


def ltn_universal_pmean(violation: torch.Tensor, p: float = 3.0) -> torch.Tensor:
    """LTN generalized p-mean universal error quantifier (Badreddine et al., AIJ 2022)."""
    return (torch.mean(violation.clamp(0.0, 1.0) ** p) + 1e-8) ** (1.0 / p)


def soft_dice_loss(logits: torch.Tensor, targets: torch.Tensor) -> torch.Tensor:
    probs = torch.sigmoid(logits)
    inter = (probs * targets).sum()
    return 1.0 - (2.0 * inter + 1.0) / (probs.sum() + targets.sum() + 1.0)


def supervised_multitask_loss(s_log: torch.Tensor, d_hat: torch.Tensor, y_s: torch.Tensor, y_d: torch.Tensor) -> torch.Tensor:
    pos_w = torch.tensor([6.0], device=s_log.device)
    bce = F.binary_cross_entropy_with_logits(s_log, y_s, pos_weight=pos_w)
    dl = soft_dice_loss(s_log, y_s)
    mse = ((1.0 + 6.0 * y_s) * (d_hat - y_d) ** 2).mean()
    return 0.45 * bce + 0.55 * dl + 2.5 * mse
\end{lstlisting}
\end{srcbox}

\section[\texorpdfstring{Unlabeled Sampling, 3D SSL Warm-Start \& Unified Training Harness (L431--L624)}{Unlabeled Sampling, 3D SSL Warm-Start & Unified Training Harness (L431--L624)}]{Unlabeled Sampling, 3D SSL Warm-Start \& Unified Training Harness (L431--L624)}
\subsection*{5.1 Test-Disjoint 3D Masked-Context SSL \& Stratified Unlabeled Bank (L431--L511)}
\addcontentsline{toc}{subsection}{\texorpdfstring{5.1 Test-Disjoint 3D Masked-Context SSL \& Stratified Unlabeled Bank (L431--L511)}{5.1 Test-Disjoint 3D Masked-Context SSL & Stratified Unlabeled Bank (L431--L511)}}
\begin{annotbox}[Leakage-Prevention \& SSL Notes --- L436--L510: Strict Held-Out Test Exclusion]
\begin{itemize}
    \item \textbf{Zero Test-Wedge Contamination (L454 \& L496):} Both \texttt{pretrain\_ssl\_backbone()} and \texttt{sample\_unlabeled\_blocks()} explicitly reject any random cube intersecting the held-out Test Wedge $[r_{x0}, r_{x1}) \times [r_{y0} + \text{split\_y}, r_{y1}) \times [r_{z0}, r_{z1})$, ensuring the test wedge remains completely unseen during both SSL pretraining and LTN semi-supervised logic training.
    \item \textbf{Biophysically Grounded SSL Targets (L441--L461, L469--L480):} Rather than reconstructing raw intensity alone, the 3-channel SSL target stacks $[I,\; (E_{\text{local}} - I)_+,\; -\text{LoG}(I)]$ (local contrast drop and 3D blobness at $\sigma = 0.48\,\text{mm}$). Inputs undergo random $35\%$ coarse $4{\times}4{\times}2$ block masking (inpainted with coarse shading channel \texttt{xb[:, 2:3]}), random gamma $[0.80, 1.25]$, and Gaussian noise ($\sigma = 0.05$), with extra loss weight on masked regions and blob structures.
    \item \textbf{Polar-to-Polar Stratification (L490--L504):} \texttt{z\_strata = [0, 8, 16, 28, 44, 60, 76, 92, 104, 112]} forces the 40-cube unlabeled LTN bank to sample both woody polar caps, tapered mid-latitudes, and the equator evenly (accepting cubes with $>12\%$ berry occupancy so sparse polar tips are included).
\end{itemize}
\end{annotbox}
\begin{srcbox}[{kiwi\_ns\_lib.py $\cdot$ Lines 431--511 --- pretrain\_ssl\_backbone() \& sample\_unlabeled\_blocks()}]
\begin{lstlisting}[style=pyannot,firstnumber=432]
# =============================================================================
# 5. Unlabeled Sampling, 3D SSL Warm-Start & Unified Training Harness
# =============================================================================

def pretrain_ssl_backbone(ds: KiwiDataset, device: torch.device, seed: int = 2027, epochs: int = 100) -> tuple[Backbone3DUNet, float]:
    """3D Masked Context Restoration SSL warm-start on 36 unlabeled test-disjoint cubes."""
    t0 = time.perf_counter()
    torch.manual_seed(seed)
    np.random.seed(seed)
    log_blob = np.clip(
        -ndi.gaussian_laplace(ds.vol_norm, sigma=(0.48 / ds.spacing[0], 0.48 / ds.spacing[1], 0.48 / ds.spacing[2])),
        0.0, 1.5,
    ).astype(np.float32)
    log_blob = np.clip(log_blob / (float(np.percentile(log_blob[ds.fruit_mask], 99.5)) + 1e-6), 0.0, 1.0)

    ssl_in, ssl_tgt = [], []
    rng = np.random.default_rng(seed)
    for _ in range(36):
        while True:
            x0 = int(rng.integers(24, 256 - 64 - 24))
            y0 = int(rng.integers(24, 256 - 40 - 24))
            z0 = int(rng.integers(16, 128 - 16 - 16))
            if (x0 < ds.rx1 and x0 + 64 > ds.rx0) and (y0 < ds.ry1 and y0 + 40 > ds.ry0 + ds.split_y) and (z0 < ds.rz1 and z0 + 16 > ds.rz0):
                continue
            if ds.fruit_mask[x0 : x0 + 64, y0 : y0 + 40, z0 : z0 + 16].mean() > 0.65:
                ssl_in.append(make_3ch_tensor(ds, x0, x0 + 64, y0, y0 + 40, z0, z0 + 16))
                c_drop = np.clip(ds.env_local[x0 : x0 + 64, y0 : y0 + 40, z0 : z0 + 16] - ds.vol_norm[x0 : x0 + 64, y0 : y0 + 40, z0 : z0 + 16], 0.0, 1.0)
                lg = log_blob[x0 : x0 + 64, y0 : y0 + 40, z0 : z0 + 16]
                ssl_tgt.append(torch.from_numpy(np.stack([ds.vol_norm[x0 : x0 + 64, y0 : y0 + 40, z0 : z0 + 16], c_drop, lg], axis=0)[None]))
                break

    x_bank = torch.cat(ssl_in, dim=0).to(device)
    y_bank = torch.cat(ssl_tgt, dim=0).to(device)
    backbone = Backbone3DUNet(in_ch=3, base_ch=16).to(device)
    head = nn.Conv3d(16, 3, kernel_size=1).to(device)
    opt = torch.optim.AdamW(list(backbone.parameters()) + list(head.parameters()), lr=4e-3, weight_decay=1e-4)

    for _ in range(epochs):
        idx = torch.randint(0, x_bank.shape[0], (4,), device=device)
        xb, yb = x_bank[idx], y_bank[idx]
        bmask = F.interpolate((torch.rand((4, 1, 16, 10, 8), device=device) < 0.35).float(), size=(64, 40, 16), mode="nearest")
        gamma = torch.empty((4, 1, 1, 1, 1), device=device).uniform_(0.80, 1.25)
        xb_cor = torch.clamp(xb.clamp_min(0.0) ** gamma + torch.randn_like(xb) * 0.05, 0.0, 1.3)
        xb_cor = xb_cor * (1.0 - bmask) + xb[:, 2:3].expand_as(xb) * bmask
        pred = head(backbone(xb_cor))
        loss = ((1.0 + 2.0 * bmask) * (1.0 + 3.0 * yb[:, 2:3]) * (pred - yb) ** 2).mean()
        opt.zero_grad()
        loss.backward()
        opt.step()

    sync_device(device)
    return backbone, time.perf_counter() - t0


def sample_unlabeled_blocks(ds: KiwiDataset, device: torch.device, n_blocks: int = 40, seed: int = 2027):
    """Sample stratified unlabeled 3D whole-fruit blocks strictly disjoint from the held-out Test Wedge."""
    ux, ug, ur, un = [], [], [], []
    rng = np.random.default_rng(seed)
    z_strata = [0, 8, 16, 28, 44, 60, 76, 92, 104, 112]
    for i in range(n_blocks):
        while True:
            x0 = int(rng.integers(8, 256 - 64 - 8))
            y0 = int(rng.integers(8, 256 - 40 - 8))
            z0 = z_strata[i % len(z_strata)]
            if (x0 < ds.rx1 and x0 + 64 > ds.rx0) and (y0 < ds.ry1 and y0 + 40 > ds.ry0 + ds.split_y) and (z0 < ds.rz1 and z0 + 16 > ds.rz0):
                continue
            if ds.fruit_mask[x0 : x0 + 64, y0 : y0 + 40, z0 : z0 + 16].mean() > 0.12:
                xa, ga = make_3ch_and_gps(ds, x0, x0 + 64, y0, y0 + 40, z0, z0 + 16)
                ux.append(xa)
                ug.append(ga)
                ur.append(torch.from_numpy(ds.soft_ring[x0 : x0 + 64, y0 : y0 + 40, z0 : z0 + 16][None, None]))
                un.append(torch.from_numpy(ds.non_dark[x0 : x0 + 64, y0 : y0 + 40, z0 : z0 + 16][None, None]))
                break
    return (
        torch.cat(ux, dim=0).to(device),
        torch.cat(ug, dim=0).to(device),
        torch.cat(ur, dim=0).to(device),
        torch.cat(un, dim=0).to(device),
    )
\end{lstlisting}
\end{srcbox}

\subsection*{5.2 Unified Multi-Mode Training Harness (L512--L624)}
\addcontentsline{toc}{subsection}{\texorpdfstring{5.2 Unified Multi-Mode Training Harness (L512--L624)}{5.2 Unified Multi-Mode Training Harness (L512--L624)}}
\begin{annotbox}[Optimization \& Experimental Control Notes --- L513--L623]
\begin{itemize}
    \item \textbf{Controlled 4-Mode Comparison (L538--L567):} Supports \texttt{"baseline"}, \texttt{"coordconv\_inorm"}, \texttt{"additive\_gps"}, and \texttt{"coupled\_gps"} under identical data splits, seeds (\texttt{torch.manual\_seed} + \texttt{np.random.seed}), \texttt{AdamW} (\texttt{weight\_decay=1e-4}), and \texttt{CosineAnnealingLR} (\texttt{eta\_min=2e-4}).
    \item \textbf{Accelerated \texttt{Organ-GPS} Learning Rate (L548--L554):} Because the 353-parameter \texttt{gps\_mlp + head\_gps} branch starts near zero while competing against a warm-started 340k-parameter CNN, giving it a separate parameter group (\texttt{gps\_lr = 1.5e-2} vs.\ \texttt{3.2e-3} backbone) lets the spatial log-prior erect crisp polar and outer-rind barriers within $\sim 25$ iterations.
    \item \textbf{Equivariant 3D Flip Augmentation (L573--L577):} Randomly flips $(X, Y, Z)$ axes with $p = 0.5$ jointly across image $\mathbf{X}_3$, coordinates $\mathbf{g}$, semantic target $y_s$, and distance target $y_d$.
    \item \textbf{Columella Hard-Negative Toggle (L582--L584, L603--L605):} Adds $0.35 \times \text{BCE}(z_{\text{col}}, 0)$ on the seed-free central columella crop; toggling \texttt{use\_columella\_neg=False} in Notebook~04 cleanly isolates spontaneous neural--symbolic load balancing.
    \item \textbf{Joint LTN Loss (L595--L611):} Samples $4$ unlabeled cubes per step and computes $\mathcal{L}_{\text{comp}} = \mathcal{L}_p((1 - R_{\text{ring}})p_u)$ and $\mathcal{L}_{T_2} = \mathcal{L}_p(p_u B_{\text{non\_dark}})$, combining $\mathcal{L}_{\text{sup}} + w_{\text{comp}}\mathcal{L}_{\text{comp}} + w_{T_2}\mathcal{L}_{T_2}$.
\end{itemize}
\end{annotbox}
\begin{srcbox}[{kiwi\_ns\_lib.py $\cdot$ Lines 512--624 --- train\_neurosymbolic\_variant()}]
\begin{lstlisting}[style=pyannot,firstnumber=513]
def train_neurosymbolic_variant(
    ds: KiwiDataset,
    device: torch.device,
    mode: str,
    ssl_backbone: Backbone3DUNet | None = None,
    unlabeled_bank: tuple[torch.Tensor, torch.Tensor, torch.Tensor, torch.Tensor] | None = None,
    p_quant: float = 3.0,
    w_comp: float = 2.0,
    w_t2: float = 1.2,
    gps_lr: float = 1.5e-2,
    epochs: int = 100,
    seed: int = 2027,
    use_columella_neg: bool = True,
) -> tuple[nn.Module, dict[str, list[float]], float]:
    """Train one controlled variant with CosineAnnealingLR and exact seed determinism."""
    torch.manual_seed(seed)
    np.random.seed(seed)
    t0 = time.perf_counter()

    train_sem_np, train_dist_np = build_distance_targets(ds.train_gt, ds.spacing)
    x_tr, g_tr = [t.to(device) for t in make_3ch_and_gps(ds, ds.rx0, ds.rx1, ds.ry0, ds.ry0 + ds.split_y, ds.rz0, ds.rz1)]
    x_ng, g_ng = [t.to(device) for t in make_3ch_and_gps(ds, ds.cx0, ds.cx1, ds.cy0, ds.cy1, ds.rz0, ds.rz1)]
    y_sem = torch.from_numpy(train_sem_np[None, None]).to(device)
    y_dist = torch.from_numpy(train_dist_np[None, None]).to(device)

    if mode == "baseline":
        model = MultiTaskUNet3D(in_ch=3, base_ch=16).to(device)
        opt = torch.optim.AdamW(model.parameters(), lr=3.2e-3, weight_decay=1e-4)
    elif mode == "coordconv_inorm":
        model = MultiTaskUNet3D(in_ch=6, base_ch=16).to(device)
        opt = torch.optim.AdamW(model.parameters(), lr=3.2e-3, weight_decay=1e-4)
    elif mode == "additive_gps":
        model = AdditiveBayesOrganNet3D(in_ch=3, gps_ch=3, base_ch=16).to(device)
        if ssl_backbone is not None:
            model.backbone.load_state_dict(ssl_backbone.state_dict())
        gps_params = list(model.gps_mlp.parameters()) + list(model.head_gps.parameters())
        gps_ids = {id(p) for p in gps_params}
        other_params = [p for p in model.parameters() if id(p) not in gps_ids]
        opt = torch.optim.AdamW(
            [{"params": other_params, "lr": 3.2e-3}, {"params": gps_params, "lr": gps_lr}],
            weight_decay=1e-4,
        )
    elif mode == "coupled_gps":
        model = CoupledOrganNet3D(in_ch=3, gps_ch=3, base_ch=16).to(device)
        if ssl_backbone is not None:
            model.backbone.load_state_dict(ssl_backbone.state_dict())
        gps_params = list(model.gps_mlp.parameters())
        gps_ids = {id(p) for p in gps_params}
        other_params = [p for p in model.parameters() if id(p) not in gps_ids]
        opt = torch.optim.AdamW(
            [{"params": other_params, "lr": 3.2e-3}, {"params": gps_params, "lr": gps_lr}],
            weight_decay=1e-4,
        )
    else:
        raise ValueError(f"Unknown mode: {mode}")

    sched = torch.optim.lr_scheduler.CosineAnnealingLR(opt, T_max=epochs, eta_min=2e-4)
    hist: dict[str, list[float]] = {"sup": [], "comp": [], "t2": []}

    for _ in range(epochs):
        dims = [d for d in (2, 3, 4) if np.random.rand() < 0.5]
        xp = torch.flip(x_tr, dims) if dims else x_tr
        gp = torch.flip(g_tr, dims) if dims else g_tr
        ys = torch.flip(y_sem, dims) if dims else y_sem
        yd = torch.flip(y_dist, dims) if dims else y_dist

        if mode == "baseline":
            sl, dh = model(xp)
            loss_sup = supervised_multitask_loss(sl, dh, ys, yd)
            if use_columella_neg:
                sl_n, _ = model(x_ng)
                loss_sup = loss_sup + 0.35 * F.binary_cross_entropy_with_logits(sl_n, torch.zeros_like(sl_n))
            loss_comp = torch.tensor(0.0, device=device)
            loss_t2 = torch.tensor(0.0, device=device)
            loss = loss_sup
        elif mode == "coordconv_inorm":
            sl, dh = model(torch.cat([xp, gp], dim=1))
            loss_sup = supervised_multitask_loss(sl, dh, ys, yd)
            if use_columella_neg:
                sl_n, _ = model(torch.cat([x_ng, g_ng], dim=1))
                loss_sup = loss_sup + 0.35 * F.binary_cross_entropy_with_logits(sl_n, torch.zeros_like(sl_n))
            ux, ug, ur, un = unlabeled_bank
            uidx = torch.randint(0, ux.shape[0], (4,), device=device)
            p_u = torch.sigmoid(model(torch.cat([ux[uidx], ug[uidx]], dim=1))[0])
            loss_comp = ltn_universal_pmean((1.0 - ur[uidx]) * p_u, p=p_quant)
            loss_t2 = ltn_universal_pmean(p_u * un[uidx], p=p_quant)
            loss = loss_sup + w_comp * loss_comp + w_t2 * loss_t2
        else:
            sl, dh = model(xp, gp)
            loss_sup = supervised_multitask_loss(sl, dh, ys, yd)
            if use_columella_neg:
                sl_n, _ = model(x_ng, g_ng)
                loss_sup = loss_sup + 0.35 * F.binary_cross_entropy_with_logits(sl_n, torch.zeros_like(sl_n))
            ux, ug, ur, un = unlabeled_bank
            uidx = torch.randint(0, ux.shape[0], (4,), device=device)
            p_u = torch.sigmoid(model(ux[uidx], ug[uidx])[0])
            loss_comp = ltn_universal_pmean((1.0 - ur[uidx]) * p_u, p=p_quant)
            loss_t2 = ltn_universal_pmean(p_u * un[uidx], p=p_quant)
            loss = loss_sup + w_comp * loss_comp + w_t2 * loss_t2

        opt.zero_grad()
        loss.backward()
        opt.step()
        sched.step()
        hist["sup"].append(float(loss_sup.item()))
        hist["comp"].append(float(loss_comp.item()))
        hist["t2"].append(float(loss_t2.item()))

    model.eval()
    sync_device(device)
    return model, hist, time.perf_counter() - t0
\end{lstlisting}
\end{srcbox}

\section[\texorpdfstring{Decoding, Declarative Hard-Core Splitting, Evaluation \& Inference (L625--L847)}{Decoding, Declarative Hard-Core Splitting, Evaluation & Inference (L625--L847)}]{Decoding, Declarative Hard-Core Splitting, Evaluation \& Inference (L625--L847)}
\subsection*{6.1 Anisotropic 3D Deep Watershed \& Declarative Mat\'ern Hard-Core Clump Splitting (L625--L690)}
\addcontentsline{toc}{subsection}{\texorpdfstring{6.1 Anisotropic 3D Deep Watershed \& Declarative Mat\'ern Hard-Core Clump Splitting (L625--L690)}{6.1 Anisotropic 3D Deep Watershed & Declarative Matern Hard-Core Clump Splitting (L625--L690)}}
\begin{annotbox}[Geometric Decoding Notes --- L630--L689: Compensating Anisotropy \& Enforcing Steric Repulsion]
\begin{itemize}
    \item \textbf{\texttt{build\_ellip()} (L630--L633):} Constructs a physical-millimeter 3D ellipsoid footprint ($r_{xy} = 0.68\,\text{mm}$, $r_z = 0.95\,\text{mm}$) so peak detection respects true millimeter dimensions despite $2.13\times$ thicker axial slices ($\Delta z = 0.50\,\text{mm}$ vs.\ $\Delta xy = 0.2344\,\text{mm}$).
    \item \textbf{\texttt{decode\_deep\_watershed\_3d()} (L636--L657):} Smooths $\hat{d} \cdot \mathbf{1}_{\text{fg}}$ lightly, detects local maxima $\ge 0.22$ within the ellipsoidal footprint, runs marker-controlled 3D watershed on $-\hat{d}_{\text{sm}}$, rescues any foreground islands $\ge 0.16\,\text{mm}^3$ that lacked a peak above $0.22$ (L642--L650), and relabels surviving instances contiguously as \texttt{uint16}.
    \item \textbf{\texttt{symbolic\_hardcore\_clump\_split()} (L660--L689):} Operationalizes Prior~3 (\texttt{Lab 5} Mat\'ern hard-core packing). Any decoded blob with $V \ge 2.55\,\text{mm}^3$ proposes $k = \operatorname{clip}(\operatorname{round}(V / 1.60), 2, 4)$ sub-clusters via $K$-means on \textbf{millimeter-scaled coordinates} (\texttt{coords * sp\_arr}). The split is accepted \emph{only} if all pairwise sub-centroid distances satisfy $d_{ab} \ge d_{\min} = 0.88\,\text{mm}$ and every sub-blob has $V_c \ge 0.35\,\text{mm}^3$---preventing over-splitting of genuinely large single seeds.
\end{itemize}
\end{annotbox}
\begin{srcbox}[{kiwi\_ns\_lib.py $\cdot$ Lines 625--690 --- build\_ellip(), decode\_deep\_watershed\_3d() \& symbolic\_hardcore\_clump\_split()}]
\begin{lstlisting}[style=pyannot,firstnumber=626]
# =============================================================================
# 6. Decoding, Declarative Hard-Core Splitting, Evaluation & Inference
# =============================================================================

def build_ellip(rx_mm: float, rz_mm: float, sp: tuple[float, float, float]) -> np.ndarray:
    rx, ry, rz = int(np.ceil(rx_mm / sp[0])), int(np.ceil(rx_mm / sp[1])), int(np.ceil(rz_mm / sp[2]))
    xx, yy, zz = np.ogrid[-rx : rx + 1, -ry : ry + 1, -rz : rz + 1]
    return ((xx * sp[0] / rx_mm) ** 2 + (yy * sp[1] / rx_mm) ** 2 + (zz * sp[2] / rz_mm) ** 2) <= 1.0


def decode_deep_watershed_3d(fg_bin: np.ndarray, dist_hat: np.ndarray, sp: tuple[float, float, float], min_mm3: float = 0.16) -> np.ndarray:
    """Marker-controlled 3D watershed seeded by anisotropic maxima of predicted distance transform."""
    d_sm = ndi.gaussian_filter(dist_hat * fg_bin, sigma=(0.65, 0.65, 0.45))
    maxima = (d_sm == ndi.maximum_filter(d_sm, footprint=build_ellip(0.68, 0.95, sp=sp))) & fg_bin & (d_sm >= 0.22)
    markers, _ = ndi.label(maxima)
    ws = segmentation.watershed(-d_sm, markers=markers, mask=fg_bin)
    unlabeled_islands, n_un = ndi.label(fg_bin & (ws == 0))
    vv = float(np.prod(sp))
    if n_un > 0:
        nxt = int(ws.max()) + 1
        for uid in range(1, n_un + 1):
            m = unlabeled_islands == uid
            if m.sum() * vv >= min_mm3:
                ws[m] = nxt
                nxt += 1
    out = np.zeros_like(ws, dtype=np.uint16)
    nid = 0
    for prop in measure.regionprops(ws):
        if prop.area * vv >= min_mm3:
            nid += 1
            out[ws == prop.label] = nid
    return out


def symbolic_hardcore_clump_split(
    lbl_3d: np.ndarray,
    sp: tuple[float, float, float],
    split_mm3: float = 2.55,
    d_min_mm: float = 0.88,
) -> np.ndarray:
    """Declarative Lab-5 Matérn hard-core clump disentanglement."""
    out = np.zeros_like(lbl_3d, dtype=np.uint16)
    sp_arr = np.asarray(sp, dtype=np.float32)
    vv = float(np.prod(sp))
    nid = 0
    for prop in measure.regionprops(lbl_3d):
        coords = prop.coords
        vol_mm3 = len(coords) * vv
        if vol_mm3 < 0.16:
            continue
        k_est = int(np.clip(round(vol_mm3 / 1.60), 2, 4)) if vol_mm3 >= split_mm3 else 1
        if k_est >= 2 and len(coords) >= k_est * 12:
            km = KMeans(n_clusters=k_est, n_init=4, random_state=2027).fit(coords.astype(np.float32) * sp_arr[None, :])
            dists = np.linalg.norm(km.cluster_centers_[:, None, :] - km.cluster_centers_[None, :, :], axis=-1)
            np.fill_diagonal(dists, 99.0)
            if dists.min() >= d_min_mm and (np.bincount(km.labels_, minlength=k_est) * vv).min() >= 0.35:
                for cid in range(k_est):
                    cp = coords[km.labels_ == cid]
                    nid += 1
                    out[cp[:, 0], cp[:, 1], cp[:, 2]] = nid
                continue
        nid += 1
        out[coords[:, 0], coords[:, 1], coords[:, 2]] = nid
    return out
\end{lstlisting}
\end{srcbox}

\subsection*{6.2 Vectorized 3D Hungarian Matching, Panoptic Quality \& Error Anatomy (L691--L751)}
\addcontentsline{toc}{subsection}{\texorpdfstring{6.2 Vectorized 3D Hungarian Matching, Panoptic Quality \& Error Anatomy (L691--L751)}{6.2 Vectorized 3D Hungarian Matching, Panoptic Quality & Error Anatomy (L691--L751)}}
\begin{annotbox}[Metric Specification --- L692--L750: Exact 3D Instance Evaluation]
\begin{itemize}
    \item \textbf{Fast Sparse Overlap Matrix (L708--L714):} Rather than looping $O(N_p \cdot N_g)$ over the 3D volume, \texttt{np.unique(np.stack([pred\_lbl[ov\_fg], gt\_lbl[ov\_fg]], axis=1), axis=0, return\_counts=True)} builds the exact intersection matrix $\mathbf{M}_{\text{inter}}$ in a single vectorized pass, then computes $\text{IoU}_{ij} = M_{ij} / (|P_i| + |G_j| - M_{ij})$.
    \item \textbf{Optimal One-to-One Assignment \& Panoptic Quality (L715--L722):} \texttt{linear\_sum\_assignment(-iou\_m)} solves the maximum-weight bipartite matching problem exactly. Matched pairs with $\text{IoU} \ge 0.25$ are counted as \texttt{TP}; then $\text{SQ} = \overline{\text{IoU}}_{\text{TP}}$, $\text{RQ} = \frac{2\,\text{TP}}{2\,\text{TP} + \text{FP} + \text{FN}}$, and $\text{PQ} = \text{SQ} \times \text{RQ}$.
    \item \textbf{Topological Error Counts (L723--L724):} A \textbf{Merge} is any single predicted blob covering $\ge 25\%$ of two or more GT seeds; a \textbf{Split} is any single GT seed overlapped $\ge 25\%$ by two or more predicted blobs.
    \item \textbf{\texttt{match\_instance\_sets()} (L733--L750):} Exposes the raw instance ID sets $(\text{TP}_{\text{pred}}, \text{TP}_{\text{gt}}, \text{FN}_{\text{gt}}, \text{FP}_{\text{pred}})$ so Notebook~04 can audit the physical volume, $T_2$ contrast drop, and $z$-border clipping of missed seeds.
\end{itemize}
\end{annotbox}
\begin{srcbox}[{kiwi\_ns\_lib.py $\cdot$ Lines 691--751 --- evaluate\_semantic\_and\_instance\_3d() \& match\_instance\_sets()}]
\begin{lstlisting}[style=pyannot,firstnumber=692]
def evaluate_semantic_and_instance_3d(pred_lbl: np.ndarray, gt_lbl: np.ndarray, min_iou: float = 0.25) -> dict:
    """Exact 3D Hungarian bipartite matching and Panoptic Quality (PQ = SQ x RQ)."""
    b_p, b_g = (pred_lbl > 0), (gt_lbl > 0)
    sem_dice = (2.0 * int((b_p & b_g).sum())) / max(int(b_p.sum() + b_g.sum()), 1)
    ids_p = [int(v) for v in np.unique(pred_lbl) if v > 0]
    ids_g = [int(v) for v in np.unique(gt_lbl) if v > 0]
    n_p, n_g = len(ids_p), len(ids_g)
    if n_p == 0 or n_g == 0:
        return {
            "Pred": n_p, "GT": n_g, "TP": 0, "FP": n_p, "FN": n_g,
            "Dice": round(sem_dice, 4), "Merges": 0, "Splits": 0,
            "Precision": 0.0, "Recall": 0.0, "SQ": 0.0, "RQ": 0.0, "PQ": 0.0,
        }
    p_map, g_map = {v: i for i, v in enumerate(ids_p)}, {v: j for j, v in enumerate(ids_g)}
    area_p = np.array([(pred_lbl == v).sum() for v in ids_p], dtype=np.float64)
    area_g = np.array([(gt_lbl == v).sum() for v in ids_g], dtype=np.float64)
    inter_m = np.zeros((n_p, n_g), dtype=np.float64)
    ov_fg = b_p & b_g
    if ov_fg.any():
        pairs, cnts = np.unique(np.stack([pred_lbl[ov_fg], gt_lbl[ov_fg]], axis=1), axis=0, return_counts=True)
        for (pv, gv), c in zip(pairs, cnts):
            inter_m[p_map[int(pv)], g_map[int(gv)]] = float(c)
    iou_m = inter_m / np.maximum(area_p[:, None] + area_g[None, :] - inter_m, 1.0)
    ri, ci = linear_sum_assignment(-iou_m)
    tp_pairs = [(r, c) for r, c in zip(ri, ci) if iou_m[r, c] >= min_iou]
    tp = len(tp_pairs)
    fp, fn = n_p - tp, n_g - tp
    prec = tp / max(tp + fp, 1)
    rec = tp / max(tp + fn, 1)
    sq = float(np.mean([iou_m[r, c] for r, c in tp_pairs])) if tp > 0 else 0.0
    rq = (2.0 * tp) / float(2 * tp + fp + fn)
    merges = int(np.sum(np.sum(inter_m / np.maximum(area_g[None, :], 1.0) >= 0.25, axis=1) >= 2))
    splits = int(np.sum(np.sum(inter_m / np.maximum(area_p[:, None], 1.0) >= 0.25, axis=0) >= 2))
    return {
        "Pred": n_p, "GT": n_g, "TP": tp, "FP": fp, "FN": fn,
        "Dice": round(sem_dice, 4), "Merges": merges, "Splits": splits,
        "Precision": round(prec, 4), "Recall": round(rec, 4),
        "SQ": round(sq, 4), "RQ": round(rq, 4), "PQ": round(sq * rq, 4),
    }


def match_instance_sets(pred_lbl: np.ndarray, gt_lbl: np.ndarray, min_iou: float = 0.25):
    """Return (tp_pred_ids, tp_gt_ids, fn_gt_ids, fp_pred_ids) for physical error anatomy."""
    ids_p = [int(v) for v in np.unique(pred_lbl) if v > 0]
    ids_g = [int(v) for v in np.unique(gt_lbl) if v > 0]
    p_map, g_map = {v: i for i, v in enumerate(ids_p)}, {v: j for j, v in enumerate(ids_g)}
    area_p = np.array([(pred_lbl == v).sum() for v in ids_p], dtype=np.float64)
    area_g = np.array([(gt_lbl == v).sum() for v in ids_g], dtype=np.float64)
    inter_m = np.zeros((len(ids_p), len(ids_g)), dtype=np.float64)
    ov = (pred_lbl > 0) & (gt_lbl > 0)
    if ov.any():
        pairs, cnts = np.unique(np.stack([pred_lbl[ov], gt_lbl[ov]], axis=1), axis=0, return_counts=True)
        for (pv, gv), c in zip(pairs, cnts):
            inter_m[p_map[int(pv)], g_map[int(gv)]] = float(c)
    iou_m = inter_m / np.maximum(area_p[:, None] + area_g[None, :] - inter_m, 1.0)
    ri, ci = linear_sum_assignment(-iou_m)
    tp_p = {ids_p[r] for r, c in zip(ri, ci) if iou_m[r, c] >= min_iou}
    tp_g = {ids_g[c] for r, c in zip(ri, ci) if iou_m[r, c] >= min_iou}
    return tp_p, tp_g, set(ids_g) - tp_g, set(ids_p) - tp_p
\end{lstlisting}
\end{srcbox}

\subsection*{6.3 Leakage-Free Operating-Point Calibration \& Unmasked 128-Slice Whole-Fruit Inference (L752--L847)}
\addcontentsline{toc}{subsection}{\texorpdfstring{6.3 Leakage-Free Operating-Point Calibration \& Unmasked 128-Slice Whole-Fruit Inference (L752--L847)}{6.3 Leakage-Free Operating-Point Calibration & Unmasked 128-Slice Whole-Fruit Inference (L752--L847)}}
\begin{annotbox}[Inference \& Validation Notes --- L753--L847: Preventing Data Leakage \& Whole-Organ Tiling]
\begin{itemize}
    \item \textbf{\texttt{forward\_patch()} (L753--L758):} Clean dispatcher hiding the signature difference between single-tensor (\texttt{baseline}, \texttt{coordconv\_inorm}) and dual-tensor (\texttt{additive\_gps}, \texttt{coupled\_gps}) models.
    \item \textbf{\texttt{calibrate\_threshold\_on\_train()} (L761--L783):} Corrects logic-induced probability calibration drift by sweeping $\tau_{\text{fg}} \in \{0.22, 0.26, \dots, 0.50\}$ \textbf{strictly on the labeled training wedge} (\texttt{ds.train\_gt}) with frozen model weights, maximizing $\text{PQ}$ without ever touching \texttt{ds.test\_gt}.
    \item \textbf{\texttt{run\_test\_and\_whole\_fruit()} (L786--L847):}
    \begin{enumerate}
        \item Evaluates strictly unmasked on the held-out Test Wedge (recording both pre-split \texttt{ev\_te\_raw} and post-split \texttt{ev\_te}), accepting \texttt{r\_override} and \texttt{pole\_override} for Notebook~05's mass-effect stress tests.
        \item When \texttt{do\_whole\_fruit=True}, tiles all 128 axial slices in $3 \times 5 \times 8 = 120$ non-overlapping $80 \times 40 \times 16$ patches ($x_0 \in \{16, 96, 176\}$, $y_0 \in \{16, 56, 96, 136, 176\}$, $z_0 \in \{0, 16, \dots, 112\}$, skipping empty air tiles with $<100$ berry voxels).
        \item Audits unmasked out-of-distribution safety directly across the whole fruit via \texttt{ood\_vox} (foreground voxels where $R_{\text{ring}} < 0.05$) and \texttt{polar\_vox} (foreground voxels in woody polar caps $z < 16$ or $z \ge 112$), then decodes full-organ 3D seed instances \texttt{ws\_wf}.
    \end{enumerate}
\end{itemize}
\end{annotbox}
\begin{warnbox}[Note on Whole-Fruit Tiling Geometry (L823--L834)]
Each inference tile ($80 \times 40 \times 16$) matches the exact shape of the training and test wedges so patch-wise \texttt{InstanceNorm3d} statistics remain consistent across the intact fruit. Notice that $x \in [16, 256)$ and $y \in [16, 216)$ span the entire berry cross-section while leaving a narrow 16-voxel air margin near the scan edges; all voxels outside \texttt{ds.fruit\_mask} are zeroed only by the intact berry boundary (\texttt{\& ds.fruit\_mask}, L836)---never by any inner-pericarp or polar mask.
\end{warnbox}
\begin{srcbox}[{kiwi\_ns\_lib.py $\cdot$ Lines 752--847 --- forward\_patch(), calibrate\_threshold\_on\_train() \& run\_test\_and\_whole\_fruit()}]
\begin{lstlisting}[style=pyannot,firstnumber=753]
def forward_patch(model: nn.Module, mode: str, x3: torch.Tensor, gps: torch.Tensor):
    if mode == "baseline":
        return model(x3)
    if mode == "coordconv_inorm":
        return model(torch.cat([x3, gps], dim=1))
    return model(x3, gps)


def calibrate_threshold_on_train(
    model: nn.Module,
    mode: str,
    ds: KiwiDataset,
    device: torch.device,
    use_split: bool = True,
) -> float:
    """Leakage-free operating-point calibration strictly on the labeled TRAIN wedge."""
    x_tr, g_tr = [t.to(device) for t in make_3ch_and_gps(ds, ds.rx0, ds.rx1, ds.ry0, ds.ry0 + ds.split_y, ds.rz0, ds.rz1)]
    with torch.no_grad():
        sl, dh = forward_patch(model, mode, x_tr, g_tr)
        p_tr = torch.sigmoid(sl)[0, 0].cpu().numpy()
        d_tr = dh[0, 0].cpu().numpy()
    best_th, best_pq = 0.42, -1.0
    for th in [0.22, 0.26, 0.30, 0.34, 0.38, 0.42, 0.46, 0.50]:
        lbl = decode_deep_watershed_3d(p_tr >= th, d_tr, ds.spacing)
        if use_split:
            lbl = symbolic_hardcore_clump_split(lbl, ds.spacing)
        ev = evaluate_semantic_and_instance_3d(lbl, ds.train_gt)
        if ev["PQ"] > best_pq:
            best_pq = ev["PQ"]
            best_th = th
    return float(best_th)


def run_test_and_whole_fruit(
    model: nn.Module,
    mode: str,
    ds: KiwiDataset,
    device: torch.device,
    thresh: float,
    use_split: bool = True,
    r_override: np.ndarray | None = None,
    pole_override: np.ndarray | None = None,
    do_whole_fruit: bool = True,
) -> dict:
    """Evaluate strictly on held-out Test Wedge (unmasked) and optionally full 128-slice berry."""
    x_te, g_te = [
        t.to(device)
        for t in make_3ch_and_gps(ds, ds.rx0, ds.rx1, ds.ry0 + ds.split_y, ds.ry1, ds.rz0, ds.rz1, r_override, pole_override)
    ]
    with torch.no_grad():
        sl_t, dh_t = forward_patch(model, mode, x_te, g_te)
        prob_te = torch.sigmoid(sl_t)[0, 0].cpu().numpy()
        dist_te = dh_t[0, 0].cpu().numpy()

    pred_te_raw = decode_deep_watershed_3d(prob_te >= thresh, dist_te, ds.spacing)
    pred_te = symbolic_hardcore_clump_split(pred_te_raw, ds.spacing) if use_split else pred_te_raw
    ev_te_raw = evaluate_semantic_and_instance_3d(pred_te_raw, ds.test_gt)
    ev_te = evaluate_semantic_and_instance_3d(pred_te, ds.test_gt)

    res = {
        "thresh": thresh, "prob_te": prob_te, "dist_te": dist_te,
        "pred_te_raw": pred_te_raw, "pred_te": pred_te,
        "ev_te_raw": ev_te_raw, "ev_te": ev_te,
    }
    if not do_whole_fruit:
        return res

    prob_wf = np.zeros(ds.vol.shape, dtype=np.float32)
    dist_wf = np.zeros(ds.vol.shape, dtype=np.float32)
    with torch.no_grad():
        for x0 in [16, 96, 176]:
            for y0 in [16, 56, 96, 136, 176]:
                for z0 in range(0, 128, 16):
                    if ds.fruit_mask[x0 : x0 + 80, y0 : y0 + 40, z0 : z0 + 16].sum() < 100:
                        continue
                    xa, ga = [
                        t.to(device)
                        for t in make_3ch_and_gps(ds, x0, x0 + 80, y0, y0 + 40, z0, z0 + 16, r_override, pole_override)
                    ]
                    s_o, d_o = forward_patch(model, mode, xa, ga)
                    prob_wf[x0 : x0 + 80, y0 : y0 + 40, z0 : z0 + 16] = torch.sigmoid(s_o)[0, 0].cpu().numpy()
                    dist_wf[x0 : x0 + 80, y0 : y0 + 40, z0 : z0 + 16] = d_o[0, 0].cpu().numpy()

    fg_wf = (prob_wf >= thresh) & ds.fruit_mask
    ood_vox = int((fg_wf & (ds.soft_ring < 0.05)).sum())
    polar_vox = int(fg_wf[:, :, :16].sum() + fg_wf[:, :, 112:].sum())
    ws_wf = decode_deep_watershed_3d(fg_wf, dist_wf, ds.spacing)
    if use_split:
        ws_wf = symbolic_hardcore_clump_split(ws_wf, ds.spacing)
    n_wf = len(np.unique(ws_wf)) - (1 if 0 in ws_wf else 0)
    res.update({
        "fg_wf": fg_wf, "ws_wf": ws_wf, "n_wf": n_wf,
        "ood_vox": ood_vox, "polar_vox": polar_vox,
    })
    return res
\end{lstlisting}
\end{srcbox}

\end{document}
